\usepackage[utf8]{inputenc}
\usepackage[T1]{fontenc}
\usepackage[portuguese]{babel}

% --- FONTE PADRÃO (Estilo Didático Encorpado) ---
\usepackage{helvet} 
\renewcommand{\familydefault}{\sfdefault}

% --- GEOMETRIA E ESPAÇAMENTO ---
\usepackage[top=1.6cm, bottom=1.5cm, left=0.9cm, right=0.9cm, headheight=22pt, headsep=0.4cm]{geometry}
\usepackage{microtype} 
\usepackage{setspace}
\setstretch{1.0}
\usepackage{parskip}
\setlength{\parskip}{2pt}
\setlength{\parindent}{0pt}

\newlength{\espacoPosTitulo}\setlength{\espacoPosTitulo}{0.3cm}
\newlength{\espacoPreQuestao}\setlength{\espacoPreQuestao}{0.1cm}
\newlength{\espacoQuestao}\setlength{\espacoQuestao}{0.2cm}
\newlength{\espacoAlt}\setlength{\espacoAlt}{0.2cm}

\usepackage{enumitem}
\setlist{itemsep=0.4cm, topsep=0.1cm, parsep=0pt, partopsep=0pt}
\newcommand{\larguraTikz}{0.85\linewidth} 
% --- TAG SEMÂNTICA PARA LEITURA DE IA (INVISÍVEL NO PDF) ---
\newenvironment{enunciadoLiteral}{}{}

% --- MATEMÁTICA E CORES ---
\usepackage{amsmath, amsfonts, amssymb, nccmath, mathastext}
\usepackage{xcolor}
\definecolor{indigo}{HTML}{4F46E5}
\definecolor{CorTitUm}{HTML}{FF6F61}
\definecolor{CorTitDois}{HTML}{FF6F61}
\definecolor{CorTitTres}{HTML}{658894}
\definecolor{CorLinha}{HTML}{B0B0B0} 
\definecolor{metal}{HTML}{7F8C8D}
\definecolor{vidro}{HTML}{3498DB}
\definecolor{ouro}{HTML}{F1C40F}
\definecolor{concreto}{HTML}{95A5A6}
\definecolor{madeira}{HTML}{D35400}
\definecolor{CinzaEscuro}{HTML}{4A4A4A} 
\definecolor{CinzaMedio}{HTML}{848484} 
\definecolor{Cinzablack}{HTML}{353535} 

% --- MINI-CABEÇALHO E RODAPÉ AUTORAL (TWOSIDE) ---
\usepackage{fancyhdr}
\pagestyle{fancy}
\fancyhf{} 
\renewcommand{\headrulewidth}{0.3pt} 
\renewcommand{\footrulewidth}{0.3pt}
\fancyhead[LE]{\small\sffamily\bfseries\color{gray!75} UNIDADE 7 - INTRODUÇÃO À ÁLGEBRA}
\ifgabarito
    \fancyhead[RO]{\small\sffamily\bfseries\color{CorTitUm}{LIVRO DO PROFESSOR -- SUPLEMENTAR}}
\else
    \fancyhead[RO]{\small\sffamily\color{gray!75} \texttt{www.profraphaelpaulino.com.br}}
\fi
\fancyfoot[LE]{\small \textbf{\thepage}}
\fancyfoot[RO]{\small \textbf{\thepage}}

% --- CAIXAS DE DESTAQUE ---
\usepackage[most]{tcolorbox}
\newtcolorbox{boxresponda}{colback=white, colframe=gray!45, boxrule=0pt, leftrule=1.7pt, sharp corners, enlarge left by=0.4cm, width=\linewidth-0.4cm, left=8pt, right=0pt, top=2pt, bottom=2pt, fontupper=\small}
\definecolor{CorDicaBorda}{HTML}{17c1be} 
\definecolor{CorDicaFundo}{HTML}{f3fbfb} 
\newtcolorbox{boxdica}{colback=CorDicaFundo, colframe=CorDicaBorda, boxrule=0pt, leftrule=2.5pt, sharp corners, width=\linewidth, left=8pt, right=8pt, top=6pt, bottom=6pt, halign=center, fontupper=\small}

% --- NOVA CAIXA PEDAGÓGICA E CORES (NOTA VERDE) ---
\definecolor{CorNotaBorda}{HTML}{10B981} 
\definecolor{CorNotaFundo}{HTML}{F0FDF4} 
\newtcolorbox{boxexplicacao}{colback=CorNotaFundo, colframe=CorNotaBorda, boxrule=0pt, leftrule=2.5pt, sharp corners, width=\linewidth, left=8pt, right=8pt, top=6pt, bottom=6pt, fontupper=\small}

% --- TÍTULOS ---
\usepackage{titlesec}
\titleformat{\section}{\color{black}\large\bfseries}{\thesection}{0em}{\vspace{0.1cm}\titlerule[0.8pt]}
\titlespacing*{\section}{0pt}{10pt}{6pt} 
\titleformat{\subsubsection}[runin]{\color{black}\bfseries}{}{0em}{}
\titlespacing*{\subsubsection}{0pt}{0pt}{0.15cm}

% --- COLUNAS E GRÁFICOS ---
\usepackage{multicol}
\setlength{\columnsep}{1cm}
\setlength{\columnseprule}{0.5pt}
\def\columnseprulecolor{\color{CorLinha}}
\usepackage{adjustbox, tikz, pgfplots, circuitikz}
\usetikzlibrary{shadows, shapes.misc, positioning, calc}
\pgfplotsset{compat=1.18}

\begin{document}
\thispagestyle{empty}
\color{Cinzablack}
\vspace*{-1.8cm}

% --- CABEÇALHO PRINCIPAL AUTORAL (Apenas 1ª página) ---
\noindent
\begin{minipage}{\textwidth}
    \fontfamily{phv}\selectfont
    \noindent
    \begin{minipage}[t]{0.65\linewidth}
        {\Large \bfseries \MakeUppercase{Caderno de Atividades Suplementares}}\\[0.1cm]
        {\large \textmd{\textcolor{CinzaEscuro}{Unidade: UNIDADE 7 - INTRODUÇÃO À ÁLGEBRA}}}
    \end{minipage}%
    \begin{minipage}[t]{0.35\linewidth}
        \raggedleft
        \ifgabarito
            {\large \bfseries \textcolor{CorTitUm}{[PROFESSOR]}}\\[0.05cm]
        \fi
        \small \textbf{Prof. Raphael Paulino}\\
        \textsf{\small \textcolor{indigo}{\texttt{profraphaelpaulino.com.br}}}
    \end{minipage}
    
    \vspace{0.15cm}
    \noindent\rule{\linewidth}{1.2pt}
    \vspace{-0.1cm}
    \footnotesize\textsf{\textbf{ÁREA:} MATEMÁTICA - ÁLGEBRA \hfill \textbf{NÍVEL:} 6º ANO}
    \vspace{0.1cm}
    \noindent\rule{\linewidth}{0.4pt}
\end{minipage}

\par\vspace{0.3cm} 
\raggedcolumns
\begin{multicols*}{2}
\vspace{0.1cm}

\subsection*{Capítulo 1: Expressões Algébricas}
\vspace{-0.2cm}\noindent\rule{\linewidth}{1.5pt} 
\par\vspace{\espacoPosTitulo}
\subsubsection*{1.} 
\begin{enunciadoLiteral}
O valor de uma expressão algébrica muda dependendo do valor numérico atribuído à sua letra (variável). Calcule o valor numérico da expressão $ 3x + 5 $ para os seguintes casos:

\begin{enumerate}[label=\alph*), itemsep=\espacoAlt]
    \item Para $ x = 2 $.
    \item Para $ x = 7 $.
\end{enumerate}
\end{enunciadoLiteral}

\ifgabarito
    % --- CAIXA LARANJA (Resolução e Tutor) ---
    \begin{tcolorbox}[colback=CorTitUm!5, colframe=CorTitUm, boxrule=1pt, arc=4pt, left=6pt, right=6pt, top=6pt, bottom=6pt]
    \small\textbf{\textcolor{CorTitUm}{Roteiro do Professor e Tutor IA:}}\par\vspace{0.1cm}
    \color{CinzaEscuro}
    \textbf{1. Ancoragem (Abordagem Socrática):} Se eu te cobrar 3 reais por cada pacote de figurinhas e somar 5 reais de frete, quanto você paga se comprar 2 pacotes? O $x$ é a quantidade de pacotes!\par\vspace{0.1cm}
    
    \textbf{2. Micro-passos da Resolução:}\\
    \textit{Passo 1:} Substituir o $x$ por 2 na expressão $\rightarrow$ $ 3 \cdot 2 + 5 $\\
    \textit{Passo 2:} Realizar a multiplicação primeiro $\rightarrow$ $ 6 + 5 $\\
    \textit{Passo 3:} Somar para obter o resultado final $\rightarrow$ $ 11 $\\
    \textit{Passo 4:} Repetir o processo para o item b, substituindo $x$ por 7 $\rightarrow$ $ 3 \cdot 7 + 5 = 21 + 5 = 26 $\par\vspace{0.1cm}
    
    \textbf{3. Mapeamento de Armadilhas (Alerta Tutor):} É muito comum o aluno no 6º ano olhar para $3x$, trocar $x$ por 2 e escrever "32" em vez de multiplicar. O tutor deve lembrar que há um sinal de multiplicação "escondido" entre o número e a letra.
    \end{tcolorbox}
    \vspace{0.15cm}
    
    % --- CAIXA VERDE (Porto Seguro Didático do Professor) ---
    \begin{boxexplicacao}
    \textbf{Porto Seguro Didático (O "Pulo do Gato"):} Na lousa, chame a letra de "caixinha mágica". Escreva $3 \cdot \text{[caixa]} + 5$. Diga que o $x$ é só uma caixa onde podemos colocar qualquer número dentro. O importante é lembrar que o número que está colado na caixa (o coeficiente 3) está multiplicando o que tem lá dentro!
    \end{boxexplicacao}
\fi

\par\vspace{\espacoQuestao}

\noindent\textcolor{CorLinha}{\rule{\linewidth}{0.5pt}} 
\par\vspace{\espacoPreQuestao}
\subsubsection*{2.} 
\begin{enunciadoLiteral}
No pátio da escola, será pintado um mural retangular. O comprimento desse mural mede o dobro da sua altura, que é representada pela letra $h$. Sendo assim, a área do mural pode ser expressa por $ 2 \cdot h \cdot h $ ou simplesmente $ 2h^2 $. Determine a área total desse mural em metros quadrados se a altura $h$ for igual a \textbf{3 m}.
\end{enunciadoLiteral}

\ifgabarito
    % --- CAIXA LARANJA (Resolução e Tutor) ---
    \begin{tcolorbox}[colback=CorTitUm!5, colframe=CorTitUm, boxrule=1pt, arc=4pt, left=6pt, right=6pt, top=6pt, bottom=6pt]
    \small\textbf{\textcolor{CorTitUm}{Roteiro do Professor e Tutor IA:}}\par\vspace{0.1cm}
    \color{CinzaEscuro}
    \textbf{1. Ancoragem (Abordagem Socrática):} Como calculamos a área de um retângulo? Se a altura é 3 e o comprimento é o dobro, qual é o comprimento?\par\vspace{0.1cm}
    
    \textbf{2. Micro-passos da Resolução:}\\
    \textit{Passo 1:} Compreender que a altura é $h=3$.\\
    \textit{Passo 2:} Substituir na expressão da área $\rightarrow$ $ 2 \cdot 3^2 $\\
    \textit{Passo 3:} Resolver a potência primeiro $\rightarrow$ $ 3^2 = 9 $\\
    \textit{Passo 4:} Multiplicar o resultado por 2 $\rightarrow$ $ 2 \cdot 9 = 18 \text{ m}^2 $\par\vspace{0.1cm}
    
    \textbf{3. Mapeamento de Armadilhas (Alerta Tutor):} O aluno pode querer multiplicar $2 \cdot 3$ primeiro e depois elevar ao quadrado, resultando em $36$. Lembre-o da ordem das operações: potências vêm antes das multiplicações.
    \end{tcolorbox}
    \vspace{0.15cm}
    
    % --- CAIXA VERDE (Porto Seguro Didático do Professor) ---
    \begin{boxexplicacao}
    \textbf{Porto Seguro Didático (O "Pulo do Gato"):} Se o aluno travar na expressão $2h^2$, faça ele desenhar! Peça para desenhar um retângulo. Coloque a altura como 3. Pergunte: "se o comprimento é o dobro, quanto é?". Ele responderá 6. Então Área = $3 \cdot 6 = 18$. Depois mostre que usar a expressão algébrica $2 \cdot 3^2$ dá exatamente o mesmo resultado de um jeito mais direto.
    \end{boxexplicacao}
\fi

\par\vspace{\espacoQuestao}

\noindent\textcolor{CorLinha}{\rule{\linewidth}{0.5pt}}
\par\vspace{\espacoPreQuestao}
\subsubsection*{3.} 
\begin{enunciadoLiteral}
Um estudante precisava calcular o valor numérico da expressão $ 4x + 7 $ para quando $ x = 5 $. Veja como ele montou a resolução no seu caderno:

\textit{"Troquei o x pelo 5. Então ficou $ 4 + 5 + 7 $. Somando tudo, cheguei ao resultado 16."}

\begin{boxresponda}
\textbf{Responda:} Qual foi a falha no conceito de Álgebra cometida por esse estudante? Corrija a resolução e apresente o resultado matematicamente correto.
\end{boxresponda}
\end{enunciadoLiteral}

\ifgabarito
    % --- CAIXA LARANJA (Resolução e Tutor) ---
    \begin{tcolorbox}[colback=CorTitUm!5, colframe=CorTitUm, boxrule=1pt, arc=4pt, left=6pt, right=6pt, top=6pt, bottom=6pt]
    \small\textbf{\textcolor{CorTitUm}{Roteiro do Professor e Tutor IA:}}\par\vspace{0.1cm}
    \color{CinzaEscuro}
    \textbf{1. Ancoragem (Abordagem Socrática):} O que significa quando um número está "grudado" em uma letra, sem nenhum sinal de mais ou de menos entre eles?\par\vspace{0.1cm}
    
    \textbf{2. Micro-passos da Resolução:}\\
    \textit{Passo 1:} Identificar o erro $\rightarrow$ O estudante achou que $4x$ significava uma soma ou apenas juntar os símbolos, escrevendo $4+5$.\\
    \textit{Passo 2:} Esclarecer a regra $\rightarrow$ O termo $4x$ significa $4 \cdot x$ (quatro vezes o valor de x).\\
    \textit{Passo 3:} Calcular corretamente $\rightarrow$ $ 4 \cdot 5 + 7 $\\
    \textit{Passo 4:} Finalizar a operação $\rightarrow$ $ 20 + 7 = 27 $\par\vspace{0.1cm}
    
    \textbf{3. Mapeamento de Armadilhas (Alerta Tutor):} Esta é uma questão metacognitiva. O aluno não pode apenas dar o número "27", ele deve obrigatoriamente escrever com palavras qual foi o erro do estudante fictício para ganhar a pontuação completa.
    \end{tcolorbox}
    \vspace{0.15cm}
    
    % --- CAIXA VERDE (Nota Pedagógica) ---
    \begin{boxexplicacao}
    \textbf{Justificativa Curricular (Raio-X da Questão):} Mede a capacidade do aluno de diagnosticar e corrigir um erro de sintaxe matemática muito comum na transição da aritmética para a álgebra (a omissão do sinal de multiplicação).
    \end{boxexplicacao}
\fi

\par\vspace{\espacoQuestao}

\noindent\textcolor{CorLinha}{\rule{\linewidth}{0.5pt}} 
\par\vspace{\espacoPreQuestao}
\subsubsection*{4.} 
\begin{enunciadoLiteral}
Em um campeonato escolar, a pontuação final de uma equipe é calculada pelo triplo do número de vitórias mais um bônus fixo de cinco pontos. Se chamarmos o número de vitórias de $v$, qual das expressões abaixo representa corretamente a pontuação final da equipe? 

Justifique por que as outras três alternativas estão incorretas, escrevendo o que elas significariam na prática.

\begin{enumerate}[label=\alph*), itemsep=\espacoAlt]
    \item $ 3(v + 5) $
    \item $ v^3 + 5 $
    \item $ 3 + v + 5 $
    \item $ 3v + 5 $
\end{enumerate}
\end{enunciadoLiteral}

\ifgabarito
    % --- CAIXA LARANJA (Resolução e Tutor) ---
    \begin{tcolorbox}[colback=CorTitUm!5, colframe=CorTitUm, boxrule=1pt, arc=4pt, left=6pt, right=6pt, top=6pt, bottom=6pt]
    \small\textbf{\textcolor{CorTitUm}{Roteiro do Professor e Tutor IA:}}\par\vspace{0.1cm}
    \color{CinzaEscuro}
    \textbf{1. Ancoragem (Abordagem Socrática):} Como você escreve "o triplo de algo" em matemática? E "mais um bônus de 5"?\par\vspace{0.1cm}
    
    \textbf{2. Micro-passos da Resolução:}\\
    \textit{Passo 1:} Validar a certa $\rightarrow$ Alternativa D. O triplo de vitórias é $3 \cdot v$ ou $3v$. Somando o bônus de 5, temos $3v + 5$.\\
    \textit{Passo 2:} Refutar A $\rightarrow$ $3(v+5)$ significa que o bônus de 5 foi somado ANTES, e a equipe ganhou o triplo de tudo (vitórias + bônus).\\
    \textit{Passo 3:} Refutar B $\rightarrow$ $v^3 + 5$ significa as vitórias elevadas ao cubo (multiplicadas por si mesmas 3 vezes), e não o triplo.\\
    \textit{Passo 4:} Refutar C $\rightarrow$ $3+v+5$ significa apenas somar 3 pontos com as vitórias, em vez de multiplicá-las por 3.\par\vspace{0.1cm}
    
    \textbf{3. Mapeamento de Armadilhas (Alerta Tutor):} O aluno pode escolher a correta (D) e esquecer de refutar as outras. O Tutor deve cobrar a justificativa escrita dos distratores, pois é aí que ocorre a verdadeira alfabetização algébrica.
    \end{tcolorbox}
    \vspace{0.15cm}
    
    % --- CAIXA VERDE (Porto Seguro Didático do Professor) ---
    \begin{boxexplicacao}
    \textbf{Porto Seguro Didático (O "Pulo do Gato"):} Na lousa, traduza a linguagem falada para o "matematiquês" pausadamente. Leia "o triplo" e escreva "3 $\cdot$ ". Leia "do número de vitórias" e escreva "$v$". Leia "mais um bônus de 5" e escreva "$+ 5$". Mostre como a frase monta a equação passo a passo.
    \end{boxexplicacao}
\fi

\par\vspace{\espacoQuestao}

\noindent\textcolor{CorLinha}{\rule{\linewidth}{0.5pt}} 
\par\vspace{\espacoPreQuestao}
\subsubsection*{5.} 
\begin{enunciadoLiteral}
A matemática se comporta como uma fábrica: entra uma matéria-prima (valor de $x$), passa por um processo (expressão algébrica) e sai um produto (resultado). 

Para a máquina cuja regra de funcionamento é a expressão algébrica $ 4x - 2 $, crie uma tabela preenchendo os resultados de saída caso os valores de entrada $x$ sejam, respectivamente, \textbf{1, 3 e 5}.
\end{enunciadoLiteral}

\ifgabarito
    % --- CAIXA LARANJA (Resolução e Tutor) ---
    \begin{tcolorbox}[colback=CorTitUm!5, colframe=CorTitUm, boxrule=1pt, arc=4pt, left=6pt, right=6pt, top=6pt, bottom=6pt]
    \small\textbf{\textcolor{CorTitUm}{Roteiro do Professor e Tutor IA:}}\par\vspace{0.1cm}
    \color{CinzaEscuro}
    \textbf{1. Ancoragem (Abordagem Socrática):} Se a máquina tem a regra "multiplique por 4 e depois tire 2", o que acontece se eu jogar o número 3 lá dentro?\par\vspace{0.1cm}
    
    \textbf{2. Micro-passos da Resolução:}\\
    \textit{Passo 1:} Para $x = 1 \rightarrow 4 \cdot 1 - 2 = 4 - 2 = 2$\\
    \textit{Passo 2:} Para $x = 3 \rightarrow 4 \cdot 3 - 2 = 12 - 2 = 10$\\
    \textit{Passo 3:} Para $x = 5 \rightarrow 4 \cdot 5 - 2 = 20 - 2 = 18$\par\vspace{0.1cm}
    
    \textbf{3. Mapeamento de Armadilhas (Alerta Tutor):} Cuidado com a ordem. A multiplicação DEVE ser feita antes da subtração. O aluno não pode fazer $1-2$ primeiro.
    
    % --- APLICANDO DROP SHADOW ---
    \begin{center}
    \begin{adjustbox}{trim=0cm 0cm 0cm 0cm, clip, max width=0.5\linewidth}
    \begin{tikzpicture}
        \definecolor{fundo}{HTML}{F8F9FA}
        \definecolor{borda}{HTML}{BDC3C7}
        \definecolor{destaque}{HTML}{2ECC71}
        
        \node[draw=borda, fill=fundo, rounded corners=3pt, drop shadow={opacity=0.15, shadow xshift=2pt, shadow yshift=-2pt}, inner sep=0pt] {
            \renewcommand{\arraystretch}{1.5}
            \begin{tabular}{c|c|c}
                \textbf{Entrada ($x$)} & \textbf{Processo ($4x-2$)} & \textbf{Saída} \\
                \hline
                $1$ & $4(1) - 2$ & \textbf{2} \\
                $3$ & $4(3) - 2$ & \textbf{10} \\
                $5$ & $4(5) - 2$ & \textbf{18} \\
            \end{tabular}
        };
    \end{tikzpicture}
    \end{adjustbox}
    \end{center}
    \end{tcolorbox}
    \vspace{0.15cm}
    
    % --- CAIXA VERDE (Porto Seguro Didático do Professor) ---
    \begin{boxexplicacao}
    \textbf{Porto Seguro Didático (O "Pulo do Gato"):} Recapitulando: A álgebra funciona como uma máquina de transformação de números. Isso prepara a intuição deles para entenderem "Funções" no 9º ano sem nenhum susto.
    \end{boxexplicacao}
\fi

\par\vspace{\espacoQuestao}

\noindent\textcolor{CorLinha}{\rule{\linewidth}{0.5pt}} 
\par\vspace{\espacoPreQuestao}
\subsubsection*{6.} 
\begin{enunciadoLiteral}
Uma empresa de engenharia utiliza blocos pré-moldados padronizados para construir pequenas passarelas. O comprimento total da passarela, em metros, é calculado somando as partes fixas das extremidades (rampas) com a parte central, que é formada por $x$ blocos encaixados. A figura ilustra o projeto.

\begin{center}
% ==========================================
% CONTROLE INDIVIDUAL DESTA IMAGEM:
% 1. CORTAR BORDAS: Mude os zeros do trim (Esquerda, Baixo, Direita, Topo). Ex: trim=1cm 0cm 1cm 0cm
% 2. TAMANHO: Para encolher só esta imagem, troque \larguraTikz por um valor (Ex: 0.5\linewidth)
% ==========================================
\begin{adjustbox}{trim=0cm 0cm 0cm 0cm, clip, max width=\larguraTikz}
\begin{tikzpicture}
    % --- APLICANDO DROP SHADOW E REGRAS PREMIUM ---
    \definecolor{concreto}{HTML}{95A5A6}
    \definecolor{asfalto}{HTML}{2F3640}
    \definecolor{grama}{HTML}{27AE60}
    
    % LAYER 1: Fundo (Grama/Chão)
    \fill[grama, rounded corners=2pt] (-2, -1.5) rectangle (10, -0.5);
    
    % LAYER 2: Passarela (Rampas fixas e blocos)
    % Rampa Esquerda (fixa 3m)
    \fill[concreto, drop shadow={opacity=0.2, shadow xshift=2pt, shadow yshift=-2pt}] (-1, -0.5) -- (0, 1) -- (2, 1) -- (2, -0.5) -- cycle;
    \draw[thick, asfalto] (-1, -0.5) -- (0, 1) -- (2, 1) -- (2, -0.5) -- cycle;
    
    % Blocos Centrais (Tamanho varia com x)
    \fill[concreto!80, drop shadow={opacity=0.2, shadow xshift=2pt, shadow yshift=-2pt}] (2.2, -0.5) rectangle (5.8, 1);
    \draw[thick, asfalto] (2.2, -0.5) rectangle (5.8, 1);
    \node[align=center, font=\bfseries\large] at (4, 0.25) {$x$ blocos\\ centrais};
    
    % Rampa Direita (fixa 2m)
    \fill[concreto, drop shadow={opacity=0.2, shadow xshift=2pt, shadow yshift=-2pt}] (6, -0.5) -- (6, 1) -- (7, 1) -- (8.5, -0.5) -- cycle;
    \draw[thick, asfalto] (6, -0.5) -- (6, 1) -- (7, 1) -- (8.5, -0.5) -- cycle;
    
    % Cotas (Medidas)
    \draw[|<->|, thick] (-1, 1.5) -- (2, 1.5) node[midway, fill=white, inner sep=2pt, rounded corners=1pt] {\textbf{3 m}};
    \draw[|<->|, thick] (2.2, 1.5) -- (5.8, 1.5) node[midway, fill=white, inner sep=2pt, rounded corners=1pt] {Cada bloco mede \textbf{4 m}};
    \draw[|<->|, thick] (6, 1.5) -- (8.5, 1.5) node[midway, fill=white, inner sep=2pt, rounded corners=1pt] {\textbf{2 m}};

\end{tikzpicture}
\end{adjustbox}
\end{center}

Com base no esquema ilustrado, se o engenheiro utilizar \textbf{5 blocos} centrais ($x=5$), qual será o comprimento total da passarela? Apresente a expressão algébrica que modela o problema e a sua resolução.
\end{enunciadoLiteral}

\ifgabarito
    % --- CAIXA LARANJA (Resolução e Tutor) ---
    \begin{tcolorbox}[colback=CorTitUm!5, colframe=CorTitUm, boxrule=1pt, arc=4pt, left=6pt, right=6pt, top=6pt, bottom=6pt]
    \small\textbf{\textcolor{CorTitUm}{Roteiro do Professor e Tutor IA:}}\par\vspace{0.1cm}
    \color{CinzaEscuro}
    \textbf{1. Ancoragem (Abordagem Socrática):} Se você olhar a imagem, o tamanho da passarela tem uma parte que nunca muda (as pontas) e uma parte que cresce se eu colocar mais blocos. Quanto mede a parte que nunca muda somada?\par\vspace{0.1cm}
    
    \textbf{2. Micro-passos da Resolução:}\\
    \textit{Passo 1:} Identificar a parte fixa visualmente $\rightarrow$ Rampa esquerda (\textbf{3 m}) + Rampa direita (\textbf{2 m}) = \textbf{5 m}.\\
    \textit{Passo 2:} Identificar a parte variável $\rightarrow$ Cada bloco mede \textbf{4 m}. Para $x$ blocos, a medida é $4 \cdot x$ ou $4x$.\\
    \textit{Passo 3:} Montar a expressão algébrica geral $\rightarrow$ Comprimento = $ 4x + 5 $.\\
    \textit{Passo 4:} Calcular para $x=5$ $\rightarrow$ $ 4 \cdot 5 + 5 = 20 + 5 = 25 \text{ metros}. $\par\vspace{0.1cm}
    
    \textbf{3. Mapeamento de Armadilhas (Alerta Tutor):} Teste Cego Absoluto aplicado. O aluno que não olhar a imagem não conseguirá resolver, pois as medidas (3m, 2m, 4m por bloco) estão exclusivamente no TikZ. Se ele errar, pergunte: "Você notou que as duas pontas da ponte têm medidas diferentes na figura?".
    \end{tcolorbox}
    \vspace{0.15cm}
    
    % --- CAIXA VERDE (Porto Seguro Didático do Professor) ---
    \begin{boxexplicacao}
    \textbf{Porto Seguro Didático (O "Pulo do Gato"):} Aqui começamos a modelagem Matemática Real. A álgebra serve justamente para isso: criar uma "fórmula" geral para um construtor não ter que pensar do zero toda vez que for montar uma ponte de tamanhos diferentes.
    \end{boxexplicacao}
\fi

\par\vspace{\espacoQuestao}

\noindent\textcolor{CorLinha}{\rule{\linewidth}{0.5pt}} 
\par\vspace{\espacoPreQuestao}
\subsubsection*{7.} 
\begin{enunciadoLiteral}
(Estilo ENEM - Adaptada) O preço de uma corrida de táxi em certa cidade é composto por uma taxa fixa, chamada de bandeirada, e um valor que depende da distância percorrida (por quilômetro). Sabendo que a bandeirada custa $ \text{R\$} \ 6,00 $ e cada quilômetro rodado custa $ \text{R\$} \ 3,00 $, o valor final $V$ da corrida, para $k$ quilômetros rodados, pode ser expresso por:
$ V = 6 + 3k $.

Dessa forma, quanto um passageiro pagará por uma viagem que percorreu uma distância de \textbf{12 quilômetros}?
\end{enunciadoLiteral}

\ifgabarito
    % --- CAIXA LARANJA (Resolução e Tutor) ---
    \begin{tcolorbox}[colback=CorTitUm!5, colframe=CorTitUm, boxrule=1pt, arc=4pt, left=6pt, right=6pt, top=6pt, bottom=6pt]
    \small\textbf{\textcolor{CorTitUm}{Roteiro do Professor e Tutor IA:}}\par\vspace{0.1cm}
    \color{CinzaEscuro}
    \textbf{1. Ancoragem (Abordagem Socrática):} Se você entrar no táxi e não andar nem um metro (0 quilômetros), você já tem que pagar alguma coisa? Por quê?\par\vspace{0.1cm}
    
    \textbf{2. Micro-passos da Resolução:}\\
    \textit{Passo 1:} Relacionar os dados do texto com a expressão dada $\rightarrow$ $k = 12$ (quilômetros rodados).\\
    \textit{Passo 2:} Substituir na expressão $\rightarrow$ $ V = 6 + 3 \cdot 12 $\\
    \textit{Passo 3:} Fazer a multiplicação do custo da distância $\rightarrow$ $ 3 \cdot 12 = 36 $\\
    \textit{Passo 4:} Somar a taxa fixa (bandeirada) $\rightarrow$ $ 6 + 36 = 42 $\\
    \textit{Passo 5:} Resposta $\rightarrow$ O passageiro pagará $\text{R\$} \ 42,00$.\par\vspace{0.1cm}
    
    \textbf{3. Mapeamento de Armadilhas (Alerta Tutor):} Alunos menos atentos leem "6 reais e 3 reais", somam logo de cara ($6+3=9$) e multiplicam por 12. Lembre-os que apenas os "3 reais" sofrem a multiplicação pela distância percorrida.
    \end{tcolorbox}
    \vspace{0.15cm}
    
    % --- CAIXA VERDE (Porto Seguro Didático do Professor) ---
    \begin{boxexplicacao}
    \textbf{Porto Seguro Didático (O "Pulo do Gato"):} Situações de Taxa Fixa + Taxa Variável são o feijão com arroz da matemática financeira básica e modelam as futuras Funções de 1º Grau. Reforce o conceito de "Fixo" (fica sozinho na conta) e "Variável" (fica multiplicando a letra).
    \end{boxexplicacao}
\fi

\par\vspace{\espacoQuestao}

\noindent\textcolor{CorLinha}{\rule{\linewidth}{0.5pt}} 
\par\vspace{\espacoPreQuestao}
\subsubsection*{8.} 
\begin{enunciadoLiteral}
Julgue as afirmativas abaixo sobre a estrutura das expressões algébricas, classificando-as como Verdadeiras (V) ou Falsas (F). Justifique matematicamente as afirmativas que você considerar falsas.

\begin{enumerate}[label=\Roman*., itemsep=\espacoAlt]
    \item Na expressão $ 5x + 10 $, o número 10 é chamado de termo constante porque seu valor não se altera, independentemente do valor de $x$.
    \item A expressão $ 2a + 3b $ significa que temos dois valores desconhecidos e iguais, representados pelas letras diferentes $a$ e $b$.
    \item Se atribuirmos $ x = 0 $ na expressão $ 8x + 4 $, o resultado da expressão será 0.
\end{enunciadoLiteral}
\end{enunciadoLiteral}

\ifgabarito
    % --- CAIXA LARANJA (Resolução e Tutor) ---
    \begin{tcolorbox}[colback=CorTitUm!5, colframe=CorTitUm, boxrule=1pt, arc=4pt, left=6pt, right=6pt, top=6pt, bottom=6pt]
    \small\textbf{\textcolor{CorTitUm}{Roteiro do Professor e Tutor IA:}}\par\vspace{0.1cm}
    \color{CinzaEscuro}
    \textbf{1. Ancoragem (Abordagem Socrática):} Se você tem caixas de dois tamanhos diferentes, você as chamaria pela mesma letra na matemática?\par\vspace{0.1cm}
    
    \textbf{2. Micro-passos da Resolução:}\\
    \textit{Passo 1:} Analisar a Afirmativa I $\rightarrow$ Verdadeira. O número solto (sem letra multiplicando) é a constante.\\
    \textit{Passo 2:} Analisar a Afirmativa II $\rightarrow$ Falsa. Letras diferentes ($a$ e $b$) representam, por definição, variáveis independentes que, em regra, guardam valores diferentes, não iguais.\\
    \textit{Passo 3:} Analisar a Afirmativa III $\rightarrow$ Falsa. Ao fazer $x = 0$, teremos $8 \cdot 0 + 4 = 0 + 4 = 4$. O resultado é 4, não zero.\par\vspace{0.1cm}
    
    \textbf{3. Mapeamento de Armadilhas (Alerta Tutor):} Na afirmativa III, o aluno pode pensar intuitivamente que "se a letra zera, tudo zera". Prove calculando para derrubar esse mito.
    \end{tcolorbox}
    \vspace{0.15cm}
    
    % --- CAIXA VERDE (Porto Seguro Didático do Professor) ---
    \begin{boxexplicacao}
    \textbf{Porto Seguro Didático (O "Pulo do Gato"):} Recapitulando: Constante é o que "fica parado" (não tem letra). Variável é o que muda. Se existem duas coisas que mudam de formas diferentes na mesma história (ex: preço de pizza e preço de refri), usamos letras diferentes para não misturar as contas.
    \end{boxexplicacao}
\fi

\par\vspace{\espacoQuestao}


\subsection*{Capítulo 2: Linguagem Algébrica}
\vspace{-0.2cm}\noindent\rule{\linewidth}{1.5pt} 
\par\vspace{\espacoPosTitulo}
\subsubsection*{9.} 
\begin{enunciadoLiteral}
O grande segredo da Álgebra é funcionar como um tradutor, transformando o Português em Matemática. Usando a letra $x$ para representar um "número desconhecido", traduza as frases abaixo para a linguagem algébrica:

\begin{enumerate}[label=\alph*), itemsep=\espacoAlt]
    \item A metade de um número mais dez.
    \item O dobro de um número, diminuído de quatro unidades.
    \item A soma de dois números consecutivos.
\end{enunciadoLiteral}
\end{enunciadoLiteral}

\ifgabarito
    % --- CAIXA LARANJA (Resolução e Tutor) ---
    \begin{tcolorbox}[colback=CorTitUm!5, colframe=CorTitUm, boxrule=1pt, arc=4pt, left=6pt, right=6pt, top=6pt, bottom=6pt]
    \small\textbf{\textcolor{CorTitUm}{Roteiro do Professor e Tutor IA:}}\par\vspace{0.1cm}
    \color{CinzaEscuro}
    \textbf{1. Ancoragem (Abordagem Socrática):} Se eu te perguntar a metade de 10, você divide por quanto? E se for a metade de $x$? O que significa ser "consecutivo" de um número que você conhece?\par\vspace{0.1cm}
    
    \textbf{2. Micro-passos da Resolução:}\\
    \textit{Passo 1:} Traduzir o item a $\rightarrow$ A metade é dividir por 2. Então: $ \mfrac{x}{2} + 10 $.\\
    \textit{Passo 2:} Traduzir o item b $\rightarrow$ O dobro é multiplicar por 2. Então: $ 2x - 4 $.\\
    \textit{Passo 3:} Traduzir o item c $\rightarrow$ Se um número é $x$, o seu consecutivo (o próximo) é $x+1$. A soma deles é $ x + (x + 1) $ ou simplificado $ 2x + 1 $.\par\vspace{0.1cm}
    
    \textbf{3. Mapeamento de Armadilhas (Alerta Tutor):} No item (c), a representação de "números consecutivos" trava muitos alunos. É vital não aceitarem $x+y$, pois $x$ e $y$ podem não ser consecutivos (ex: 2 e 9). Para garantir que são seguidos, usa-se a regra de $+1$.
    \end{tcolorbox}
    \vspace{0.15cm}
    
    % --- CAIXA VERDE (Porto Seguro Didático do Professor) ---
    \begin{boxexplicacao}
    \textbf{Porto Seguro Didático (O "Pulo do Gato"):} Ensine o aluno a ler "picotado". Mostre a frase: [A metade de um número] | [mais] | [dez]. Cada bloco vira uma parte matemática. O consecutivo merece um exemplo real: "Se eu tenho o 5, o consecutivo é 6 (que é 5+1). Se eu tenho o x, o consecutivo é x+1".
    \end{boxexplicacao}
\fi

\par\vspace{\espacoQuestao}

\noindent\textcolor{CorLinha}{\rule{\linewidth}{0.5pt}} 
\par\vspace{\espacoPreQuestao}
\subsubsection*{10.} 
\begin{enunciadoLiteral}
A figura a seguir ilustra a planta baixa de um canteiro retangular em uma praça pública. As medidas das laterais estão representadas por expressões algébricas, em metros.

\begin{center}
% ==========================================
% CONTROLE INDIVIDUAL DESTA IMAGEM:
% ==========================================
\begin{adjustbox}{trim=0cm 0cm 0cm 0cm, clip, max width=0.6\linewidth}
\begin{tikzpicture}
    % --- APLICANDO DROP SHADOW E REGRAS PREMIUM ---
    \definecolor{cimento}{HTML}{BDC3C7}
    \definecolor{grama}{HTML}{2ECC71}
    
    % LAYER 1: Canteiro com borda arredondada
    \fill[grama, rounded corners=3pt, drop shadow={opacity=0.2, shadow xshift=3pt, shadow yshift=-3pt}] (0,0) rectangle (6,3);
    \draw[thick, cimento, rounded corners=3pt] (0,0) rectangle (6,3);
    
    % LAYER 2: Textos das medidas, usando máscara de fundo para leitura
    \node[fill=white, inner sep=2pt, rounded corners=2pt, font=\bfseries] at (3, -0.4) {$2x + 1$};
    \node[fill=white, inner sep=2pt, rounded corners=2pt, font=\bfseries] at (6.6, 1.5) {$x$};
    \node[fill=white, inner sep=2pt, rounded corners=2pt, font=\bfseries] at (3, 3.4) {$2x + 1$};
    \node[fill=white, inner sep=2pt, rounded corners=2pt, font=\bfseries] at (-0.6, 1.5) {$x$};
    
    % Cotas ilustrativas (linhas finas)
    \draw[dashed, thin, gray] (0,-0.1) -- (0,-0.4) -- (2.3,-0.4);
    \draw[dashed, thin, gray] (6,-0.1) -- (6,-0.4) -- (3.7,-0.4);
    \draw[dashed, thin, gray] (6.1,0) -- (6.6,0) -- (6.6,1.2);
    \draw[dashed, thin, gray] (6.1,3) -- (6.6,3) -- (6.6,1.8);
\end{tikzpicture}
\end{adjustbox}
\end{center}

Determine a expressão algébrica simplificada que representa o \textbf{perímetro} total deste canteiro. Se for necessário, agrupe os termos semelhantes (junte todas as letras $x$ e todos os números sozinhos).
\end{enunciadoLiteral}

\ifgabarito
    % --- CAIXA LARANJA (Resolução e Tutor) ---
    \begin{tcolorbox}[colback=CorTitUm!5, colframe=CorTitUm, boxrule=1pt, arc=4pt, left=6pt, right=6pt, top=6pt, bottom=6pt]
    \small\textbf{\textcolor{CorTitUm}{Roteiro do Professor e Tutor IA:}}\par\vspace{0.1cm}
    \color{CinzaEscuro}
    \textbf{1. Ancoragem (Abordagem Socrática):} O que é perímetro na geometria? Se eu for dar uma volta completa andando pela borda desse jardim, quais são as 4 medidas que vou somar?\par\vspace{0.1cm}
    
    \textbf{2. Micro-passos da Resolução:}\\
    \textit{Passo 1:} Definir perímetro $\rightarrow$ É a soma de todos os 4 lados do retângulo.\\
    \textit{Passo 2:} Escrever a soma completa baseada na imagem $\rightarrow$ Perímetro = $ (2x + 1) + (x) + (2x + 1) + (x) $\\
    \textit{Passo 3:} Agrupar os termos semelhantes (família do x com família do x, números com números) $\rightarrow$ $ 2x + x + 2x + x $ e $ 1 + 1 $\\
    \textit{Passo 4:} Reduzir a expressão $\rightarrow$ Perímetro = $ 6x + 2 $.\par\vspace{0.1cm}
    
    \textbf{3. Mapeamento de Armadilhas (Alerta Tutor):} O aluno pode confundir perímetro com área e tentar multiplicar as expressões. Foque firmemente no conceito de "contorno". Outro erro comum é somar apenas a base e uma altura ($3x+1$) esquecendo os outros dois lados paralelos.
    \end{tcolorbox}
    \vspace{0.15cm}
    
    % --- CAIXA VERDE (Porto Seguro Didático do Professor) ---
    \begin{boxexplicacao}
    \textbf{Porto Seguro Didático (O "Pulo do Gato"):} "Termos semelhantes" é como separar frutas. "x" são maçãs, números sozinhos são laranjas. Você não pode somar 2 maçãs com 1 laranja e dizer que tem "3 maçã-laranjas". Some todos os $x$ de um lado ($2x+1x+2x+1x = 6x$) e os números do outro ($1+1=2$).
    \end{boxexplicacao}
\fi

\par\vspace{\espacoQuestao}

\noindent\textcolor{CorLinha}{\rule{\linewidth}{0.5pt}}
\par\vspace{\espacoPreQuestao}
\subsubsection*{11.} 
\begin{enunciadoLiteral}
Uma professora desenhou um quadrado na lousa e indicou que a medida do seu lado vale $3a$. Em seguida, ela pediu para a turma escrever a expressão que representa a \textbf{área} dessa figura.

O aluno Joãozinho rapidamente respondeu: "A área do quadrado é $ 3a \cdot 4 $, o que resulta em $ 12a $."

\begin{boxresponda}
\textbf{Responda:} Joãozinho acertou? Caso ele tenha cometido um equívoco conceitual entre área e perímetro, explique qual foi a sua confusão e apresente a expressão algébrica correta para a Área deste quadrado.
\end{boxresponda}
\end{enunciadoLiteral}

\ifgabarito
    % --- CAIXA LARANJA (Resolução e Tutor) ---
    \begin{tcolorbox}[colback=CorTitUm!5, colframe=CorTitUm, boxrule=1pt, arc=4pt, left=6pt, right=6pt, top=6pt, bottom=6pt]
    \small\textbf{\textcolor{CorTitUm}{Roteiro do Professor e Tutor IA:}}\par\vspace{0.1cm}
    \color{CinzaEscuro}
    \textbf{1. Ancoragem (Abordagem Socrática):} Qual é a diferença na fórmula entre o "contorno" de um quadrado e a "superfície interna" dele?\par\vspace{0.1cm}
    
    \textbf{2. Micro-passos da Resolução:}\\
    \textit{Passo 1:} Diagnosticar o erro $\rightarrow$ Joãozinho não acertou. Ele calculou o Perímetro (soma dos 4 lados ou lado vezes 4), mas a professora pediu a Área.\\
    \textit{Passo 2:} Lembrar a fórmula da Área $\rightarrow$ A área de um quadrado é lado $\cdot$ lado (ou lado ao quadrado).\\
    \textit{Passo 3:} Construir a conta correta $\rightarrow$ Área = $ 3a \cdot 3a $\\
    \textit{Passo 4:} Multiplicar corretamente (número com número, letra com letra) $\rightarrow$ $ 3 \cdot 3 = 9 $ e $ a \cdot a = a^2 $.\\
    \textit{Passo 5:} Resposta final $\rightarrow$ A expressão correta da Área é $ 9a^2 $.\par\vspace{0.1cm}
    
    \textbf{3. Mapeamento de Armadilhas (Alerta Tutor):} No $6^\circ$ ano, potências com letras ($a \cdot a = a^2$) causam certa hesitação. É importante mostrar que, assim como $5 \cdot 5 = 5^2$, $a \cdot a = a^2$. O Joãozinho fez $3a \cdot 4$ (perímetro).
    \end{tcolorbox}
    \vspace{0.15cm}
    
    % --- CAIXA VERDE (Nota Pedagógica) ---
    \begin{boxexplicacao}
    \textbf{Justificativa Curricular (Raio-X da Questão):} Avalia de forma simultânea a geometria plana (diferenciação de grandezas de 1D vs 2D) e a operação algébrica básica com monômios.
    \end{boxexplicacao}
\fi

\par\vspace{\espacoQuestao}

\noindent\textcolor{CorLinha}{\rule{\linewidth}{0.5pt}} 
\par\vspace{\espacoPreQuestao}
\subsubsection*{12.} 
\begin{enunciadoLiteral}
Considere um terreno retangular onde um dos lados mede \textbf{5 metros} e o outro lado mede $ (x + 2) \textbf{ metros} $. Qual das alternativas abaixo apresenta a expressão algébrica correta da Área desse terreno, caso apliquemos a propriedade distributiva ("chuveirinho")?

Prove por que a alternativa escolhida está certa refutando, por meio de cálculos, as falhas das alternativas incorretas.

\begin{enumerate}[label=\alph*), itemsep=\espacoAlt]
    \item $ 5x + 2 $
    \item $ x + 10 $
    \item $ 5x + 10 $
    \item $ 5 + x + 2 $
\end{enumerate}
\end{enunciadoLiteral}

\ifgabarito
    % --- CAIXA LARANJA (Resolução e Tutor) ---
    \begin{tcolorbox}[colback=CorTitUm!5, colframe=CorTitUm, boxrule=1pt, arc=4pt, left=6pt, right=6pt, top=6pt, bottom=6pt]
    \small\textbf{\textcolor{CorTitUm}{Roteiro do Professor e Tutor IA:}}\par\vspace{0.1cm}
    \color{CinzaEscuro}
    \textbf{1. Ancoragem (Abordagem Socrática):} Para achar a área de um retângulo multiplicamos a base pela altura. Se eu for multiplicar o 5 por toda a expressão $(x+2)$, o 5 multiplica só o $x$, ou multiplica o 2 também?\par\vspace{0.1cm}
    
    \textbf{2. Micro-passos da Resolução:}\\
    \textit{Passo 1:} Montar o cálculo da Área $\rightarrow$ $ 5 \cdot (x + 2) $.\\
    \textit{Passo 2:} Validar a certa (C) $\rightarrow$ Aplicando a propriedade distributiva (o "chuveirinho"): $5 \cdot x + 5 \cdot 2 = 5x + 10$.\\
    \textit{Passo 3:} Refutar A $\rightarrow$ Em $5x + 2$, a pessoa multiplicou o 5 apenas pelo $x$ e "esqueceu" de aplicar a distributiva no número 2.\\
    \textit{Passo 4:} Refutar B $\rightarrow$ Em $x + 10$, a pessoa multiplicou o 5 apenas pelo 2, esquecendo do $x$.\\
    \textit{Passo 5:} Refutar D $\rightarrow$ Em $5 + x + 2$, a pessoa confundiu as regras geométricas e apenas tentou somar os termos em vez de multiplicá-los.\par\vspace{0.1cm}
    
    \textbf{3. Mapeamento de Armadilhas (Alerta Tutor):} A falha de distributiva parcial (distrair-se e multiplicar apenas o primeiro termo do parêntese, gerando a alternativa A) é o erro algébrico mais duradouro até o Ensino Médio. Corte o mal pela raiz exigindo o desenho das "setinhas" indicando as multiplicações.
    \end{tcolorbox}
    \vspace{0.15cm}
    
    % --- CAIXA VERDE (Porto Seguro Didático do Professor) ---
    \begin{boxexplicacao}
    \textbf{Porto Seguro Didático (O "Pulo do Gato"):} Use a metáfora da entrega de correio. O 5, que está fora do parêntese, é o carteiro. Os termos $x$ e $2$, que estão dentro da casa (parênteses), são moradores. O carteiro tem que entregar a carta para TODOS os moradores da casa. Se ele entregar só para o $x$ (gerando $5x+2$), o morador "2" vai ficar bravo!
    \end{boxexplicacao}
\fi

\par\vspace{\espacoQuestao}
\subsubsection*{13.} 
\begin{enunciadoLiteral}
Uma fábrica embala sucos de laranja em caixas de papelão. Cada caixa fechada contém exatamente $x$ garrafas. Para atender ao pedido de um supermercado, a fábrica enviou um lote composto por \textbf{5 caixas} completas e mais \textbf{3 garrafas} avulsas. 

Qual expressão algébrica representa o total de garrafas enviadas neste pedido? 
\end{enunciadoLiteral}

\ifgabarito
    % --- CAIXA LARANJA (Resolução e Tutor) ---
    \begin{tcolorbox}[colback=CorTitUm!5, colframe=CorTitUm, boxrule=1pt, arc=4pt, left=6pt, right=6pt, top=6pt, bottom=6pt]
    \small\textbf{\textcolor{CorTitUm}{Roteiro do Professor e Tutor IA:}}\par\vspace{0.1cm}
    \color{CinzaEscuro}
    \textbf{1. Ancoragem (Abordagem Socrática):} Se eu te der 5 caixas fechadas e dentro de cada uma tiver, digamos, 10 garrafas, que conta você faz para descobrir o total nas caixas? E se em vez de 10 for $x$?\par\vspace{0.1cm}
    
    \textbf{2. Micro-passos da Resolução:}\\
    \textit{Passo 1:} Quantificar o que está dentro das caixas $\rightarrow$ São 5 caixas, cada uma com $x$ garrafas. Isso significa $ 5 \cdot x $ ou $ 5x $.\\
    \textit{Passo 2:} Adicionar as garrafas soltas $\rightarrow$ Somamos 3 ao valor das caixas.\\
    \textit{Passo 3:} Escrever a expressão final $\rightarrow$ Total = $ 5x + 3 $.\par\vspace{0.1cm}
    
    \textbf{3. Mapeamento de Armadilhas (Alerta Tutor):} Alunos tendem a querer somar tudo, gerando $5+3+x$ ou $8x$. Reforce a ideia de "grupos" (multiplicação) versus "unidades avulsas" (adição).
    \end{tcolorbox}
    \vspace{0.15cm}
    
    % --- CAIXA VERDE (Porto Seguro Didático do Professor) ---
    \begin{boxexplicacao}
    \textbf{Porto Seguro Didático (O "Pulo do Gato"):} O uso de "caixas" e "itens avulsos" é a melhor imagem mental para $ax + b$. A caixa esconde o valor de $x$ (não sabemos quanto tem lá dentro), e o número que acompanha o $x$ é quantas caixas temos. O número solto (o 3) é o que ficou fora da caixa e a gente consegue contar a olho nu.
    \end{boxexplicacao}
\fi

\par\vspace{\espacoQuestao}

\noindent\textcolor{CorLinha}{\rule{\linewidth}{0.5pt}} 
\par\vspace{\espacoPreQuestao}
\subsubsection*{14.} 
\begin{enunciadoLiteral}
(Estilo VUNESP) Para realizar uma viagem escolar, o custo do aluguel do ônibus foi dividido igualmente entre os alunos participantes. O motorista cobrou uma taxa fixa de $ \text{R\$} \ 250,00 $ pelo serviço e mais $ \text{R\$} \ 15,00 $ por cada aluno que entrasse no ônibus. Sabendo que $x$ representa o número de alunos que foram à viagem, qual das alternativas apresenta a expressão algébrica do valor \textbf{total} cobrado pelo motorista?

Justifique a resposta correta e demonstre matematicamente o erro das alternativas falsas.

\begin{enumerate}[label=\alph*), itemsep=\espacoAlt]
    \item $ 250 \cdot 15x $
    \item $ 250 + 15 + x $
    \item $ 15x + 250 $
    \item $ 250x + 15 $
\end{enumerate}
\end{enunciadoLiteral}

\ifgabarito
    % --- CAIXA LARANJA (Resolução e Tutor) ---
    \begin{tcolorbox}[colback=CorTitUm!5, colframe=CorTitUm, boxrule=1pt, arc=4pt, left=6pt, right=6pt, top=6pt, bottom=6pt]
    \small\textbf{\textcolor{CorTitUm}{Roteiro do Professor e Tutor IA:}}\par\vspace{0.1cm}
    \color{CinzaEscuro}
    \textbf{1. Ancoragem (Abordagem Socrática):} Se fossem apenas 2 alunos, você pagaria 15 reais uma vez só, ou duas vezes? E os 250 reais, mudam se forem 2 ou 10 alunos?\par\vspace{0.1cm}
    
    \textbf{2. Micro-passos da Resolução:}\\
    \textit{Passo 1:} Validar a certa $\rightarrow$ Alternativa C. O valor de 15 reais multiplica a quantidade de alunos ($15x$) e soma-se à taxa fixa de 250, gerando $15x + 250$.\\
    \textit{Passo 2:} Refutar A $\rightarrow$ A alternativa multiplica a taxa fixa pela variável ($250 \cdot 15x$), o que seria um absurdo econômico (uma taxa fixa multiplicando um valor que já multiplica os alunos).\\
    \textit{Passo 3:} Refutar B $\rightarrow$ Ela propõe que o preço por aluno não multiplique a quantidade de alunos, mas apenas some $15 + x$. Se fossem 2 alunos, seriam apenas 17 reais em vez de $15 \cdot 2$.\\
    \textit{Passo 4:} Refutar D $\rightarrow$ Aqui a taxa fixa e a variável foram invertidas. Estão cobrando $250$ por aluno e fixando $15$ reais, o oposto do texto.\par\vspace{0.1cm}
    
    \textbf{3. Mapeamento de Armadilhas (Alerta Tutor):} Inverter o coeficiente variável com o termo constante (Alternativa D) é um erro de leitura muito comum. O aluno lê 250 primeiro no texto e gruda o $x$ nele instintivamente. Peça atenção à palavra "por cada".
    \end{tcolorbox}
    \vspace{0.15cm}
    
    % --- CAIXA VERDE (Porto Seguro Didático do Professor) ---
    \begin{boxexplicacao}
    \textbf{Porto Seguro Didático (O "Pulo do Gato"):} Recapitulando a aula anterior: a palavra mágica para descobrir quem ganha o "$x$" é "POR" ou "CADA". 15 reais *por cada* aluno significa que o 15 será multiplicado ($15x$). O valor que vem escrito como "taxa fixa" fica sempre sozinho, sem letra.
    \end{boxexplicacao}
\fi

\par\vspace{\espacoQuestao}

\noindent\textcolor{CorLinha}{\rule{\linewidth}{0.5pt}} 
\par\vspace{\espacoPreQuestao}
\subsubsection*{15.} 
\begin{enunciadoLiteral}
Um maratonista treina correndo trajetos diferentes na cidade. O mapa abaixo mostra o percurso do corredor, saindo do ponto \textbf{A}, passando pelo ponto \textbf{B} e chegando ao ponto \textbf{C}. A distância entre os pontos está indicada por expressões algébricas em quilômetros.

\begin{center}
% ==========================================
% CONTROLE INDIVIDUAL DESTA IMAGEM:
% ==========================================
\begin{adjustbox}{trim=0cm 0cm 0cm 0cm, clip, max width=0.7\linewidth}
\begin{tikzpicture}
    % --- APLICANDO DROP SHADOW E REGRAS PREMIUM ---
    \definecolor{asfalto}{HTML}{34495E}
    \definecolor{ponto}{HTML}{E74C3C}
    \definecolor{trilha}{HTML}{F1C40F}
    
    % LAYER 1: Caminho (Rua)
    \draw[line width=10pt, asfalto, drop shadow={opacity=0.2, shadow xshift=2pt, shadow yshift=-2pt}] (0,0) -- (4,2) -- (8,0);
    \draw[line width=2pt, dashed, trilha] (0,0) -- (4,2) -- (8,0);
    
    % LAYER 2: Pontos A, B, C
    \filldraw[ponto, draw=white, thick] (0,0) circle (6pt) node[above left, black, font=\bfseries] {Ponto A};
    \filldraw[ponto, draw=white, thick] (4,2) circle (6pt) node[above, black, font=\bfseries, yshift=2pt] {Ponto B};
    \filldraw[ponto, draw=white, thick] (8,0) circle (6pt) node[above right, black, font=\bfseries] {Ponto C};
    
    % LAYER 3: Textos das medidas (Cotas)
    \node[fill=white, inner sep=3pt, rounded corners=2pt, font=\bfseries] at (1.8, 1.4) {$2x + 1$};
    \node[fill=white, inner sep=3pt, rounded corners=2pt, font=\bfseries] at (6.2, 1.4) {$3x - 1$};
    
\end{tikzpicture}
\end{adjustbox}
\end{center}

Se a distância total percorrida pelo maratonista do Ponto A até o Ponto C for de \textbf{15 quilômetros}, escreva a igualdade (equação) matemática que descreve essa situação, traduzindo o mapa para a Álgebra. \textit{(Atenção: Não é necessário resolver a equação agora, apenas monte a igualdade)}
\end{enunciadoLiteral}

\ifgabarito
    % --- CAIXA LARANJA (Resolução e Tutor) ---
    \begin{tcolorbox}[colback=CorTitUm!5, colframe=CorTitUm, boxrule=1pt, arc=4pt, left=6pt, right=6pt, top=6pt, bottom=6pt]
    \small\textbf{\textcolor{CorTitUm}{Roteiro do Professor e Tutor IA:}}\par\vspace{0.1cm}
    \color{CinzaEscuro}
    \textbf{1. Ancoragem (Abordagem Socrática):} Se eu te disser que caminhei um pedaço e depois caminhei outro pedaço, como eu descubro o total? O que o símbolo de "igual" significa em um problema onde eu já te dou o resultado final?\par\vspace{0.1cm}
    
    \textbf{2. Micro-passos da Resolução:}\\
    \textit{Passo 1:} Identificar a primeira parte do trajeto (A até B) $\rightarrow$ $ 2x + 1 $.\\
    \textit{Passo 2:} Identificar a segunda parte do trajeto (B até C) $\rightarrow$ $ 3x - 1 $.\\
    \textit{Passo 3:} Montar a soma dos trajetos $\rightarrow$ $ (2x + 1) + (3x - 1) $.\\
    \textit{Passo 4:} Igualar ao total dado no texto para montar a equação $\rightarrow$ $ (2x + 1) + (3x - 1) = 15 $.\par\vspace{0.1cm}
    
    \textbf{3. Mapeamento de Armadilhas (Alerta Tutor):} Os alunos no início da álgebra esquecem de colocar o sinal de "$=$" e o total ("$= 15$"). Eles montam apenas a expressão da soma. Lembre-os de que a diferença entre uma *Expressão* e uma *Equação* é a presença de uma balança de igualdade.
    \end{tcolorbox}
    \vspace{0.15cm}
    
    % --- CAIXA VERDE (Porto Seguro Didático do Professor) ---
    \begin{boxexplicacao}
    \textbf{Porto Seguro Didático (O "Pulo do Gato"):} Esta questão é a ponte dourada entre Expressões e Equações. Mostre que até agora nós brincávamos de montar a conta sem saber o fim dela. Na Equação, nós SABEMOS o fim (é 15). O trabalho do detetive agora será descobrir quem é o $x$ que faz essa soma mágica dar 15.
    \end{boxexplicacao}
\fi

\par\vspace{\espacoQuestao}


\subsection*{Capítulo 3: Equações do 1º Grau - Introdução}
\vspace{-0.2cm}\noindent\rule{\linewidth}{1.5pt} 
\par\vspace{\espacoPosTitulo}
\subsubsection*{16.} 
\begin{enunciadoLiteral}
Uma equação funciona exatamente como uma balança de pratos que está em perfeito equilíbrio. Observe a balança a seguir: de um lado temos duas caixas de massa desconhecida $x$ e um peso de \textbf{5 kg}. Do outro lado, temos um único peso de \textbf{17 kg}.

\begin{center}
% ==========================================
% CONTROLE INDIVIDUAL DESTA IMAGEM:
% ==========================================
\begin{adjustbox}{trim=0cm 0cm 0cm 0cm, clip, max width=0.65\linewidth}
\begin{tikzpicture}
    % --- APLICANDO DROP SHADOW E REGRAS PREMIUM ---
    \definecolor{madeira}{HTML}{D35400}
    \definecolor{metal}{HTML}{7F8C8D}
    \definecolor{peso}{HTML}{34495E}
    \definecolor{caixa}{HTML}{F1C40F}
    
    % LAYER 1: Base e fulcro da balança
    \fill[metal, drop shadow={opacity=0.3, shadow xshift=2pt, shadow yshift=-2pt}] (-1, -0.5) rectangle (1, 0);
    \fill[metal, drop shadow={opacity=0.2, shadow xshift=2pt, shadow yshift=-2pt}] (-0.2, 0) -- (0.2, 0) -- (0, 2) -- cycle;
    
    % LAYER 2: Haste e pratos (Equilíbrio horizontal)
    \draw[line width=4pt, metal] (-4, 2) -- (4, 2);
    % Cordas prato esquerdo
    \draw[thick, gray] (-4, 2) -- (-5, 0.5);
    \draw[thick, gray] (-4, 2) -- (-3, 0.5);
    % Cordas prato direito
    \draw[thick, gray] (4, 2) -- (3, 0.5);
    \draw[thick, gray] (4, 2) -- (5, 0.5);
    
    % Pratos
    \fill[madeira, rounded corners=2pt] (-5.5, 0.3) rectangle (-2.5, 0.5);
    \fill[madeira, rounded corners=2pt] (2.5, 0.3) rectangle (5.5, 0.5);
    
    % LAYER 3: Pesos no prato ESQUERDO
    \fill[caixa, rounded corners=2pt, drop shadow={opacity=0.2, shadow xshift=1pt, shadow yshift=-1pt}] (-4.8, 0.5) rectangle (-3.8, 1.5);
    \node[font=\bfseries\large] at (-4.3, 1) {$x$};
    
    \fill[caixa, rounded corners=2pt, drop shadow={opacity=0.2, shadow xshift=1pt, shadow yshift=-1pt}] (-3.7, 0.5) rectangle (-2.7, 1.5);
    \node[font=\bfseries\large] at (-3.2, 1) {$x$};
    
    \fill[peso, rounded corners=1pt] (-5.4, 0.5) -- (-4.9, 0.5) -- (-5.0, 1) -- (-5.3, 1) -- cycle;
    \node[font=\bfseries, text=white, scale=0.7] at (-5.15, 0.75) {5kg};
    
    % LAYER 4: Pesos no prato DIREITO
    \fill[peso, rounded corners=2pt, drop shadow={opacity=0.2, shadow xshift=2pt, shadow yshift=-2pt}] (3.5, 0.5) -- (4.5, 0.5) -- (4.3, 1.8) -- (3.7, 1.8) -- cycle;
    \node[font=\bfseries, text=white, scale=1.1] at (4, 1.1) {17kg};
    
\end{tikzpicture}
\end{adjustbox}
\end{center}

Escreva a equação matemática que representa o equilíbrio dessa balança e, usando a lógica das operações (retirar pesos iguais dos dois lados), descubra a massa $x$ de cada caixa.
\end{enunciadoLiteral}

\ifgabarito
    % --- CAIXA LARANJA (Resolução e Tutor) ---
    \begin{tcolorbox}[colback=CorTitUm!5, colframe=CorTitUm, boxrule=1pt, arc=4pt, left=6pt, right=6pt, top=6pt, bottom=6pt]
    \small\textbf{\textcolor{CorTitUm}{Roteiro do Professor e Tutor IA:}}\par\vspace{0.1cm}
    \color{CinzaEscuro}
    \textbf{1. Ancoragem (Abordagem Socrática):} Se eu tirar o peso de 5 kg do prato da esquerda, a balança vai desequilibrar. O que eu tenho que fazer no prato da direita para manter o equilíbrio exato?\par\vspace{0.1cm}
    
    \textbf{2. Micro-passos da Resolução:}\\
    \textit{Passo 1:} Montar a equação visual $\rightarrow$ O prato esquerdo é $x + x + 5$ ou $2x + 5$. O direito é $17$. A igualdade é: $ 2x + 5 = 17 $.\\
    \textit{Passo 2:} Usar a lógica do equilíbrio (Princípio Aditivo) $\rightarrow$ Retiramos 5 kg de AMBOS os lados para isolar as caixas: $ 2x = 17 - 5 $.\\
    \textit{Passo 3:} Calcular a diferença $\rightarrow$ $ 2x = 12 $. (Duas caixas pesam 12 kg).\\
    \textit{Passo 4:} Usar o Princípio Multiplicativo $\rightarrow$ Se duas caixas iguais pesam 12, para achar uma, dividimos por 2: $ x = \mfrac{12}{2} $.\\
    \textit{Passo 5:} Resposta $\rightarrow$ $ x = 6 \text{ kg} $.\par\vspace{0.1cm}
    
    \textbf{3. Mapeamento de Armadilhas (Alerta Tutor):} No $6^\circ$ ano, a regra do "passa para o outro lado mudando o sinal" ainda NÃO DEVE ser ensinada de forma decorada. A justificativa rigorosa deve ser sempre a manutenção da balança: "Tudo o que eu faço de um lado, faço do outro".
    \end{tcolorbox}
    \vspace{0.15cm}
    
    % --- CAIXA VERDE (Porto Seguro Didático do Professor) ---
    \begin{boxexplicacao}
    \textbf{Porto Seguro Didático (O "Pulo do Gato"):} Desenhe na lousa uma balança simples. O objetivo de toda equação é deixar o $x$ SOZINHO em um dos pratos. Como? Limpando o terreno. Primeiro, tira os números soltos. Tirou de um lado? Obrigatoriamente tira do outro (senão a balança cai). Depois de tirar os "intrusos", é só dividir!
    \end{boxexplicacao}
\fi

\par\vspace{\espacoQuestao}

\noindent\textcolor{CorLinha}{\rule{\linewidth}{0.5pt}} 
\par\vspace{\espacoPreQuestao}
\subsubsection*{17.} 
\begin{enunciadoLiteral}
Em matemática, a "raiz" (ou solução) de uma equação é o valor exato que devemos colocar no lugar do $x$ para que a balança de igualdade seja verdadeira. 

Verifique, realizando os cálculos, se o número \textbf{5} é a raiz da equação $ 3x + 2 = 17 $.
\end{enunciadoLiteral}

\ifgabarito
    % --- CAIXA LARANJA (Resolução e Tutor) ---
    \begin{tcolorbox}[colback=CorTitUm!5, colframe=CorTitUm, boxrule=1pt, arc=4pt, left=6pt, right=6pt, top=6pt, bottom=6pt]
    \small\textbf{\textcolor{CorTitUm}{Roteiro do Professor e Tutor IA:}}\par\vspace{0.1cm}
    \color{CinzaEscuro}
    \textbf{1. Ancoragem (Abordagem Socrática):} Em vez de tentar "resolver" a conta, vamos fazer o caminho inverso. O que acontece se a gente simplesmente pegar o 5, colocar na "caixinha" do $x$ e calcular? O resultado dá 17?\par\vspace{0.1cm}
    
    \textbf{2. Micro-passos da Resolução:}\\
    \textit{Passo 1:} Substituir $x$ por 5 no primeiro membro da equação $\rightarrow$ $ 3 \cdot 5 + 2 $.\\
    \textit{Passo 2:} Realizar a multiplicação primeiro $\rightarrow$ $ 15 + 2 $.\\
    \textit{Passo 3:} Realizar a soma $\rightarrow$ $ 17 $.\\
    \textit{Passo 4:} Comparar com o segundo membro da equação original $\rightarrow$ $ 17 = 17 $. É uma igualdade verdadeira.\\
    \textit{Passo 5:} Conclusão $\rightarrow$ Sim, o número 5 é a raiz da equação.\par\vspace{0.1cm}
    
    \textbf{3. Mapeamento de Armadilhas (Alerta Tutor):} O aluno pode tentar resolver a equação isolando o $x$, errar na conta e concluir falsamente que 5 não é a raiz. A orientação deve ser estrita: problemas de "verificar" não se resolvem de trás pra frente, apenas substituímos o número sugerido e validamos a igualdade (prova real).
    \end{tcolorbox}
    \vspace{0.15cm}
    
    % --- CAIXA VERDE (Porto Seguro Didático do Professor) ---
    \begin{boxexplicacao}
    \textbf{Porto Seguro Didático (O "Pulo do Gato"):} "Verificar a raiz" é igual a testar uma chave na fechadura. Você não precisa desmontar a porta para saber se a chave funciona, basta colocá-la na maçaneta (substituir o x) e ver se a porta abre (se a igualdade final dá lados iguais).
    \end{boxexplicacao}
\fi

\par\vspace{\espacoQuestao}

\noindent\textcolor{CorLinha}{\rule{\linewidth}{0.5pt}}
\par\vspace{\espacoPreQuestao}
\subsubsection*{18.} 
\begin{enunciadoLiteral}
Um estudante estava tentando resolver a equação $ x + 8 = 20 $. Para tentar deixar a letra $x$ isolada no lado esquerdo da balança, ele fez a seguinte anotação no caderno:

\textit{"Eu preciso tirar o 8 que está atrapalhando o $x$. Então, eu retiro 8 do lado esquerdo e, para compensar, eu somo 8 do lado direito, chegando na resposta $ x = 28 $."}

\begin{boxresponda}
\textbf{Responda:} Baseando-se no princípio da balança em equilíbrio, qual foi o erro de raciocínio lógico cometido por este estudante? Qual é o valor real de $x$?
\end{boxresponda}
\end{enunciadoLiteral}

\ifgabarito
    % --- CAIXA LARANJA (Resolução e Tutor) ---
    \begin{tcolorbox}[colback=CorTitUm!5, colframe=CorTitUm, boxrule=1pt, arc=4pt, left=6pt, right=6pt, top=6pt, bottom=6pt]
    \small\textbf{\textcolor{CorTitUm}{Roteiro do Professor e Tutor IA:}}\par\vspace{0.1cm}
    \color{CinzaEscuro}
    \textbf{1. Ancoragem (Abordagem Socrática):} Imagine uma gangorra perfeitamente equilibrada com duas pessoas de pesos iguais. Se você tira uma mochila de 8 kg das costas da pessoa da esquerda e PENDURA nas costas da pessoa da direita, a gangorra vai continuar em equilíbrio?\par\vspace{0.1cm}
    
    \textbf{2. Micro-passos da Resolução:}\\
    \textit{Passo 1:} Diagnosticar o erro $\rightarrow$ O estudante feriu o Princípio de Equivalência. Ele subtraiu de um lado e SOMOU do outro. Isso desequilibra completamente a balança.\\
    \textit{Passo 2:} Explicar a regra correta $\rightarrow$ O que se faz de um lado, deve-se fazer \textbf{exatamente igual} do outro. Se tirou 8 da esquerda, TEM que tirar 8 da direita.\\
    \textit{Passo 3:} Resolver corretamente $\rightarrow$ $ x + 8 - 8 = 20 - 8 $.\\
    \textit{Passo 4:} Resposta final $\rightarrow$ $ x = 12 $.\par\vspace{0.1cm}
    
    \textbf{3. Mapeamento de Armadilhas (Alerta Tutor):} Este é o clássico reflexo condicionado precoce do "passar mudando o sinal" ensinado de forma errada nos anos anteriores, onde a criança confunde as operações. Esta questão metacognitiva quebra o "macete" e resgata o conceito real.
    \end{tcolorbox}
    \vspace{0.15cm}
    
    % --- CAIXA VERDE (Nota Pedagógica) ---
    \begin{boxexplicacao}
    \textbf{Justificativa Curricular (Raio-X da Questão):} Uma habilidade vital no 6º ano é descontruir algoritmos cegos (os macetes). O aluno que entende o porquê de diminuir 8 de ambos os lados raramente cometerá erros de sinal mais à frente.
    \end{boxexplicacao}
\fi

\par\vspace{\espacoQuestao}

\noindent\textcolor{CorLinha}{\rule{\linewidth}{0.5pt}} 
\par\vspace{\espacoPreQuestao}
\subsubsection*{19.} 
\begin{enunciadoLiteral}
Tente descobrir mentalmente (ou fazendo os cálculos reversos) a solução do seguinte desafio algébrico: "O quádruplo de um número misterioso, subtraído de 2 unidades, resulta em 18". 

Ou seja, $ 4x - 2 = 18 $.

Qual das opções abaixo corresponde ao número misterioso $x$? Justifique provando matematicamente que as outras opções NÃO funcionam.

\begin{enumerate}[label=\alph*), itemsep=\espacoAlt]
    \item $ x = 3 $
    \item $ x = 4 $
    \item $ x = 5 $
    \item $ x = 6 $
\end{enumerate}
\end{enunciadoLiteral}

\ifgabarito
    % --- CAIXA LARANJA (Resolução e Tutor) ---
    \begin{tcolorbox}[colback=CorTitUm!5, colframe=CorTitUm, boxrule=1pt, arc=4pt, left=6pt, right=6pt, top=6pt, bottom=6pt]
    \small\textbf{\textcolor{CorTitUm}{Roteiro do Professor e Tutor IA:}}\par\vspace{0.1cm}
    \color{CinzaEscuro}
    \textbf{1. Ancoragem (Abordagem Socrática):} Que número que multiplicado por 4 e tirando 2 dá 18? Vamos testar os candidatos?\par\vspace{0.1cm}
    
    \textbf{2. Micro-passos da Resolução:}\\
    \textit{Passo 1:} Testar e refutar A ($x=3$) $\rightarrow$ $ 4 \cdot 3 - 2 = 12 - 2 = 10 $ (Falso, precisava dar 18).\\
    \textit{Passo 2:} Testar e refutar B ($x=4$) $\rightarrow$ $ 4 \cdot 4 - 2 = 16 - 2 = 14 $ (Falso).\\
    \textit{Passo 3:} Validar a alternativa correta C ($x=5$) $\rightarrow$ $ 4 \cdot 5 - 2 = 20 - 2 = 18 $ (Verdadeiro, balança equilibrou).\\
    \textit{Passo 4:} Testar e refutar D ($x=6$) $\rightarrow$ $ 4 \cdot 6 - 2 = 24 - 2 = 22 $ (Falso).\par\vspace{0.1cm}
    
    \textbf{3. Mapeamento de Armadilhas (Alerta Tutor):} O aluno que tentar resolver pela lógica operacional invertida e falhar ($18-2=16 \rightarrow x=4$) deve ser convidado a fazer o teste de substituição. Ele próprio perceberá a falha da sua lógica ao substituir o 4 e dar 14.
    \end{tcolorbox}
    \vspace{0.15cm}
    
    % --- CAIXA VERDE (Porto Seguro Didático do Professor) ---
    \begin{boxexplicacao}
    \textbf{Porto Seguro Didático (O "Pulo do Gato"):} "Método da Tentativa e Erro Estruturado". No início da Álgebra, forçar o aluno a chutar valores organizadamente cria nele o "senso de proporção". Ele percebe que ao chutar 3, deu um número baixo. Ao chutar 4, aproximou-se. Assim, ele intui que a equação tem um comportamento linear!
    \end{boxexplicacao}
\fi

\par\vspace{\espacoQuestao}

\noindent\textcolor{CorLinha}{\rule{\linewidth}{0.5pt}} 
\par\vspace{\espacoPreQuestao}
\subsubsection*{20.} 
\begin{enunciadoLiteral}
Uma fita retangular foi dividida em partes. A primeira parte tem tamanho de \textbf{15 cm}. As outras três partes menores têm tamanhos iguais, representados por $x$. Sabendo que a fita inteira possui \textbf{30 cm} no total, observe o esquema em blocos:

\begin{center}
% ==========================================
% CONTROLE INDIVIDUAL DESTA IMAGEM:
% ==========================================
\begin{adjustbox}{trim=0cm 0cm 0cm 0cm, clip, max width=0.7\linewidth}
\begin{tikzpicture}
    % --- APLICANDO DROP SHADOW E REGRAS PREMIUM ---
    \definecolor{fitaG}{HTML}{3498DB}
    \definecolor{fitaP}{HTML}{E67E22}
    
    % Fita inteira (fundo comparativo)
    \fill[gray!20, rounded corners=2pt, drop shadow={opacity=0.1, shadow xshift=2pt, shadow yshift=-2pt}] (0, 1.5) rectangle (8, 2.5);
    \draw[thick, gray!50, rounded corners=2pt] (0, 1.5) rectangle (8, 2.5);
    \node[font=\bfseries, text=black] at (4, 2) {Total = 30 cm};
    
    % Peças recortadas
    \fill[fitaG, rounded corners=2pt, drop shadow={opacity=0.3, shadow xshift=2pt, shadow yshift=-2pt}] (0, 0) rectangle (4, 1);
    \draw[thick, black] (0, 0) rectangle (4, 1);
    \node[font=\bfseries, text=white] at (2, 0.5) {15 cm};
    
    \fill[fitaP, rounded corners=2pt, drop shadow={opacity=0.3, shadow xshift=2pt, shadow yshift=-2pt}] (4.2, 0) rectangle (5.3, 1);
    \draw[thick, black] (4.2, 0) rectangle (5.3, 1);
    \node[font=\bfseries, text=white] at (4.75, 0.5) {$x$};
    
    \fill[fitaP, rounded corners=2pt, drop shadow={opacity=0.3, shadow xshift=2pt, shadow yshift=-2pt}] (5.5, 0) rectangle (6.6, 1);
    \draw[thick, black] (5.5, 0) rectangle (6.6, 1);
    \node[font=\bfseries, text=white] at (6.05, 0.5) {$x$};
    
    \fill[fitaP, rounded corners=2pt, drop shadow={opacity=0.3, shadow xshift=2pt, shadow yshift=-2pt}] (6.8, 0) rectangle (7.9, 1);
    \draw[thick, black] (6.8, 0) rectangle (7.9, 1);
    \node[font=\bfseries, text=white] at (7.35, 0.5) {$x$};
    
    % Chave de agrupamento (brace)
    \draw [decorate, decoration={brace, amplitude=5pt}, thick] (7.9, -0.2) -- (0, -0.2) node [midway, below, yshift=-4pt, font=\bfseries] {Comprimento equivalente};
\end{tikzpicture}
\end{adjustbox}
\end{center}

Com base na figura, escreva a equação correspondente e determine o valor (tamanho) de cada pedaço $x$.
\end{enunciadoLiteral}

\ifgabarito
    % --- CAIXA LARANJA (Resolução e Tutor) ---
    \begin{tcolorbox}[colback=CorTitUm!5, colframe=CorTitUm, boxrule=1pt, arc=4pt, left=6pt, right=6pt, top=6pt, bottom=6pt]
    \small\textbf{\textcolor{CorTitUm}{Roteiro do Professor e Tutor IA:}}\par\vspace{0.1cm}
    \color{CinzaEscuro}
    \textbf{1. Ancoragem (Abordagem Socrática):} Se a fita toda tem 30 cm e a parte azul gasta 15 cm, quantos centímetros "sobram" para serem divididos pelos blocos laranjas?\par\vspace{0.1cm}
    
    \textbf{2. Micro-passos da Resolução:}\\
    \textit{Passo 1:} Montar a equação lendo a imagem $\rightarrow$ Temos três blocos de $x$ somados a um de $15$, que juntos formam o total $30$. Equação: $ 3x + 15 = 30 $.\\
    \textit{Passo 2:} Isolar as incógnitas $\rightarrow$ Subtraímos 15 dos dois lados (ou "o pedaço grande de 15 gasta metade da fita, sobram 15 para o resto"): $ 3x = 30 - 15 \Rightarrow 3x = 15 $.\\
    \textit{Passo 3:} Dividir os bloquinhos $\rightarrow$ Se três blocos valem 15, dividimos 15 por 3: $ x = \mfrac{15}{3} $.\\
    \textit{Passo 4:} Resposta $\rightarrow$ $ x = 5 \text{ cm} $.\par\vspace{0.1cm}
    
    \textbf{3. Mapeamento de Armadilhas (Alerta Tutor):} A modelagem visual (Método de Singapura) ajuda muito, mas o aluno pode tentar adivinhar a resposta de cabeça sem escrever a equação. Exija a escrita formal: $3x + 15 = 30$. O objetivo não é apenas achar 5, é alfabetizar algebricamente.
    \end{tcolorbox}
    \vspace{0.15cm}
    
    % --- CAIXA VERDE (Porto Seguro Didático do Professor) ---
    \begin{boxexplicacao}
    \textbf{Porto Seguro Didático (O "Pulo do Gato"):} Este formato em barras retangulares é mundialmente famoso como "Matemática de Singapura". Ele materializa a álgebra, mostrando que resolver uma equação é apenas "recortar" o pedaço que você conhece (15) do total (30), e o que sobra você "fatia" igualmente (por 3).
    \end{boxexplicacao}
\fi

\par\vspace{\espacoQuestao}

\noindent\textcolor{CorLinha}{\rule{\linewidth}{0.5pt}} 
\par\vspace{\espacoPreQuestao}
\subsubsection*{21.} 
\begin{enunciadoLiteral}
Em uma brincadeira matemática de adivinhação, Marina disse a seguinte frase ao seu irmão: 

\textit{"Pensei em um número natural. Multipliquei esse número por dois e, em seguida, adicionei seis. O resultado final dessa conta foi vinte."}

Transforme o desafio de Marina em uma equação do 1º grau (chamando o número pensado de $x$) e resolva para descobrir em qual número ela pensou.
\end{enunciadoLiteral}

\ifgabarito
    % --- CAIXA LARANJA (Resolução e Tutor) ---
    \begin{tcolorbox}[colback=CorTitUm!5, colframe=CorTitUm, boxrule=1pt, arc=4pt, left=6pt, right=6pt, top=6pt, bottom=6pt]
    \small\textbf{\textcolor{CorTitUm}{Roteiro do Professor e Tutor IA:}}\par\vspace{0.1cm}
    \color{CinzaEscuro}
    \textbf{1. Ancoragem (Abordagem Socrática):} Se ela no final somou 6 para chegar no 20, que número ela tinha *antes* de somar esse 6?\par\vspace{0.1cm}
    
    \textbf{2. Micro-passos da Resolução:}\\
    \textit{Passo 1:} Tradução (Montar a Equação) $\rightarrow$ O dobro do número ($2x$) mais seis ($+ 6$) é igual a vinte ($= 20$). Equação: $ 2x + 6 = 20 $.\\
    \textit{Passo 2:} Isolamento do termo com letra $\rightarrow$ Subtraindo 6 dos dois lados da balança (ou passando na operação inversa): $ 2x = 20 - 6 $.\\
    \textit{Passo 3:} Calcular a diferença $\rightarrow$ $ 2x = 14 $.\\
    \textit{Passo 4:} Dividir $\rightarrow$ O número que multiplicado por 2 dá 14 é $ 14 \div 2 $. Logo, $ x = 7 $.\\
    \textit{Passo 5:} Resposta $\rightarrow$ Marina pensou no número 7.\par\vspace{0.1cm}
    
    \textbf{3. Mapeamento de Armadilhas (Alerta Tutor):} O aluno tentará fazer a "operação inversa" de cabeça e pode errar a ordem, dividindo o 20 por 2 primeiro e depois tirando 6 (chegando em 4). Alerte que a ordem de desfazer é inversa à ordem de fazer: se o $+6$ foi o último passo, ele é o primeiro a ser desfeito.
    \end{tcolorbox}
    \vspace{0.15cm}
    
    % --- CAIXA VERDE (Porto Seguro Didático do Professor) ---
    \begin{boxexplicacao}
    \textbf{Porto Seguro Didático (O "Pulo do Gato"):} Resolver equações é como "desfazer as malas" depois de uma viagem. Você tira por último o que colocou primeiro. Na brincadeira, a Marina primeiro multiplicou, e por último somou 6. Na hora de resolver, nós tiramos o 6 primeiro, e por último desfazemos a multiplicação!
    \end{boxexplicacao}
\fi

\par\vspace{\espacoQuestao}

\noindent\textcolor{CorLinha}{\rule{\linewidth}{0.5pt}} 
\par\vspace{\espacoPreQuestao}
\subsubsection*{22.} 
\begin{enunciadoLiteral}
(Estilo FUVEST - Adaptada) Um cliente comprou um caderno por $ \text{R\$} \ 15,00 $ e \textbf{quatro canetas idênticas} de mesmo valor, gastando um total de $ \text{R\$} \ 35,00 $. Se representarmos o valor de cada caneta por $y$, qual das equações abaixo representa corretamente essa compra e qual o valor de cada caneta? 

Demonstre por meio de cálculos o motivo de as outras três equações estarem erradas em relação à história narrada.

\begin{enumerate}[label=\alph*), itemsep=\espacoAlt]
    \item $ 15y + 4 = 35 $
    \item $ 4y = 35 + 15 $
    \item $ 4y + 15 = 35 $
    \item $ y + 15 = 35 $
\end{enumerate}
\end{enunciadoLiteral}

\ifgabarito
    % --- CAIXA LARANJA (Resolução e Tutor) ---
    \begin{tcolorbox}[colback=CorTitUm!5, colframe=CorTitUm, boxrule=1pt, arc=4pt, left=6pt, right=6pt, top=6pt, bottom=6pt]
    \small\textbf{\textcolor{CorTitUm}{Roteiro do Professor e Tutor IA:}}\par\vspace{0.1cm}
    \color{CinzaEscuro}
    \textbf{1. Ancoragem (Abordagem Socrática):} Se o caderno já custou 15, quanto sobrou do seu dinheiro (dos 35) apenas para pagar as quatro canetas?\par\vspace{0.1cm}
    
    \textbf{2. Micro-passos da Resolução:}\\
    \textit{Passo 1:} Validar a certa (C) $\rightarrow$ 4 canetas vezes $y$ mais 15 reais é igual a 35 reais ($4y + 15 = 35$).\\
    \textit{Passo 2:} Resolver a C $\rightarrow$ $ 4y = 35 - 15 \Rightarrow 4y = 20 \Rightarrow y = 5 $. Cada caneta custou 5 reais.\\
    \textit{Passo 3:} Refutar A $\rightarrow$ Em $15y + 4 = 35$, a pessoa está comprando 15 canetas e somando 4 reais de taxa (inverteu as quantidades).\\
    \textit{Passo 4:} Refutar B $\rightarrow$ Em $4y = 35 + 15$, estaríamos dizendo que as canetas sozinhas custaram 50 reais (a soma da compra com o caderno), o que é impossível pois o total era 35.\\
    \textit{Passo 5:} Refutar D $\rightarrow$ Em $y + 15 = 35$, a pessoa comprou apenas UMA caneta, esquecendo-se do coeficiente 4.\par\vspace{0.1cm}
    
    \textbf{3. Mapeamento de Armadilhas (Alerta Tutor):} Cuidado extremo com a alternativa B. Visualmente ela atrai alunos que lembram da regra de "passar o número pra lá", mas eles erram o sinal, somando o 15 em vez de subtraí-lo do total de 35.
    \end{tcolorbox}
    \vspace{0.15cm}
    
    % --- CAIXA VERDE (Porto Seguro Didático do Professor) ---
    \begin{boxexplicacao}
    \textbf{Porto Seguro Didático (O "Pulo do Gato"):} O erro da alternativa B ($35+15$) é fantástico para mostrar por que não podemos apenas decorar regras de "sinal". Se você comprou tudo por 35 reais, como é que só as canetas vão custar 50? O bom senso matemático é o melhor fiscal da álgebra.
    \end{boxexplicacao}
\fi

\par\vspace{\espacoQuestao}

\noindent\textcolor{CorLinha}{\rule{\linewidth}{0.5pt}} 
\par\vspace{\espacoPreQuestao}
\subsubsection*{23.} 
\begin{enunciadoLiteral}
Um carpinteiro construiu uma moldura de madeira no formato de um triângulo, onde a medida de um lado é desconhecida e representada por $x$. Os outros dois lados medem, em metros, o dobro do primeiro lado ($2x$) e o triplo do primeiro lado ($3x$), conforme ilustrado no desenho a seguir.

\begin{center}
% ==========================================
% CONTROLE INDIVIDUAL DESTA IMAGEM:
% ==========================================
\begin{adjustbox}{trim=0cm 0cm 0cm 0cm, clip, max width=0.6\linewidth}
\begin{tikzpicture}
    % --- APLICANDO DROP SHADOW E REGRAS PREMIUM ---
    \definecolor{madeiraTri}{HTML}{D35400}
    \definecolor{madeiraBorda}{HTML}{A04000}
    
    % LAYER 1: Triângulo moldura
    \filldraw[fill=madeiraTri!80, draw=madeiraBorda, line width=4pt, line join=round, drop shadow={opacity=0.3, shadow xshift=3pt, shadow yshift=-2pt}] (0,0) -- (6,0) -- (1.5,4) -- cycle;
    
    % LAYER 2: Textos vazados
    \node[fill=white, inner sep=2pt, rounded corners=2pt, font=\bfseries\large] at (3, -0.4) {$3x$};
    \node[fill=white, inner sep=2pt, rounded corners=2pt, font=\bfseries\large, rotate=68] at (0.35, 2.3) {$x$};
    \node[fill=white, inner sep=2pt, rounded corners=2pt, font=\bfseries\large, rotate=-40] at (4.2, 2.5) {$2x$};

\end{tikzpicture}
\end{adjustbox}
\end{center}

Se o carpinteiro utilizou exatamente \textbf{12 metros} lineares de madeira para construir a volta completa dessa moldura (ou seja, seu perímetro), monte a equação adequada e descubra a medida, em metros, do menor lado ($x$).
\end{enunciadoLiteral}

\ifgabarito
    % --- CAIXA LARANJA (Resolução e Tutor) ---
    \begin{tcolorbox}[colback=CorTitUm!5, colframe=CorTitUm, boxrule=1pt, arc=4pt, left=6pt, right=6pt, top=6pt, bottom=6pt]
    \small\textbf{\textcolor{CorTitUm}{Roteiro do Professor e Tutor IA:}}\par\vspace{0.1cm}
    \color{CinzaEscuro}
    \textbf{1. Ancoragem (Abordagem Socrática):} Perímetro é a soma do contorno. Quantos "x" nós temos somando os três lados juntos da figura?\par\vspace{0.1cm}
    
    \textbf{2. Micro-passos da Resolução:}\\
    \textit{Passo 1:} Montar a soma dos lados $\rightarrow$ $ x + 2x + 3x $.\\
    \textit{Passo 2:} Igualar ao perímetro total dado no texto $\rightarrow$ $ x + 2x + 3x = 12 $.\\
    \textit{Passo 3:} Reduzir os termos semelhantes (somar a família do x) $\rightarrow$ Lembre-se que $x$ sozinho é $1x$. Assim, $1x + 2x + 3x = 6x$. Logo, $ 6x = 12 $.\\
    \textit{Passo 4:} Isolar a variável $\rightarrow$ $ x = \mfrac{12}{6} $.\\
    \textit{Passo 5:} Resposta final $\rightarrow$ O menor lado ($x$) mede \textbf{2 metros}.\par\vspace{0.1cm}
    
    \textbf{3. Mapeamento de Armadilhas (Alerta Tutor):} O erro mais letal aqui é o aluno somar $2x + 3x$ dando $5x$, ignorando olimpicamente que a letra $x$ sozinha carrega um $1$ oculto na frente. Se o aluno encontrar $5x=12$ (não resultando em natural), questione: "E aquele primeiro lado do triângulo, ele não vale nada?".
    \end{tcolorbox}
    \vspace{0.15cm}
    
    % --- CAIXA VERDE (Porto Seguro Didático do Professor) ---
    \begin{boxexplicacao}
    \textbf{Porto Seguro Didático (O "Pulo do Gato"):} O Fantasma do Número 1! A álgebra tem fantasmas amigáveis. Na frente de toda letra $x$ sozinha, existe um fantasma de jaleco escrito $1$. Portanto, $x+2x$ nunca dá $2x$, dá $3x$. É só somar $1$ lousa mais $2$ lousas.
    \end{boxexplicacao}
\fi

\par\vspace{\espacoQuestao}


\subsection*{Capítulo 4: Sequências Numéricas}
\vspace{-0.2cm}\noindent\rule{\linewidth}{1.5pt} 
\par\vspace{\espacoPosTitulo}
\subsubsection*{24.} 
\begin{enunciadoLiteral}
Em uma aula de lógica, os alunos formaram figuras montando quadrados com palitos de fósforo. Observe a lei de formação desta sequência visual:

\begin{center}
% ==========================================
% CONTROLE INDIVIDUAL DESTA IMAGEM:
% ==========================================
\begin{adjustbox}{trim=0cm 0cm 0cm 0cm, clip, max width=0.8\linewidth}
\begin{tikzpicture}
    % --- APLICANDO DROP SHADOW E REGRAS PREMIUM ---
    \definecolor{palito}{HTML}{E74C3C}
    \definecolor{cabeca}{HTML}{8E44AD}
    \tikzset{
        stick/.style={line width=2.5pt, palito, drop shadow={opacity=0.2, shadow xshift=1pt, shadow yshift=-1pt}},
        dot/.style={circle, fill=cabeca, inner sep=1.5pt}
    }
    
    % FIGURA 1 (1 quadrado = 4 palitos)
    \node[font=\bfseries, text=black] at (0.5, 1.5) {Figura 1};
    \draw[stick] (0,0) -- (1,0); \node[dot] at (1,0) {};
    \draw[stick] (1,0) -- (1,1); \node[dot] at (1,1) {};
    \draw[stick] (1,1) -- (0,1); \node[dot] at (0,1) {};
    \draw[stick] (0,1) -- (0,0); \node[dot] at (0,0) {};
    
    % FIGURA 2 (2 quadrados = 7 palitos)
    \node[font=\bfseries, text=black] at (3.5, 1.5) {Figura 2};
    % quad 1
    \draw[stick] (2.5,0) -- (3.5,0); \node[dot] at (3.5,0) {};
    \draw[stick] (3.5,0) -- (3.5,1); \node[dot] at (3.5,1) {};
    \draw[stick] (3.5,1) -- (2.5,1); \node[dot] at (2.5,1) {};
    \draw[stick] (2.5,1) -- (2.5,0); \node[dot] at (2.5,0) {};
    % quad 2 adicionado
    \draw[stick] (3.5,0) -- (4.5,0); \node[dot] at (4.5,0) {};
    \draw[stick] (4.5,0) -- (4.5,1); \node[dot] at (4.5,1) {};
    \draw[stick] (4.5,1) -- (3.5,1); 
    
    % FIGURA 3 (3 quadrados = 10 palitos)
    \node[font=\bfseries, text=black] at (7.5, 1.5) {Figura 3};
    % quad 1
    \draw[stick] (6,0) -- (7,0); \node[dot] at (7,0) {};
    \draw[stick] (7,0) -- (7,1); \node[dot] at (7,1) {};
    \draw[stick] (7,1) -- (6,1); \node[dot] at (6,1) {};
    \draw[stick] (6,1) -- (6,0); \node[dot] at (6,0) {};
    % quad 2
    \draw[stick] (7,0) -- (8,0); \node[dot] at (8,0) {};
    \draw[stick] (8,0) -- (8,1); \node[dot] at (8,1) {};
    \draw[stick] (8,1) -- (7,1); 
    % quad 3 adicionado
    \draw[stick] (8,0) -- (9,0); \node[dot] at (9,0) {};
    \draw[stick] (9,0) -- (9,1); \node[dot] at (9,1) {};
    \draw[stick] (9,1) -- (8,1); 
    
\end{tikzpicture}
\end{adjustbox}
\end{center}

Analisando o padrão de crescimento, para adicionar um novo quadrado adjacente à sequência, quantos palitos "novos" são acrescentados na estrutura? 

Descubra quantos palitos no total seriam necessários para montar a \textbf{Figura 5} usando essa mesma regularidade numérica.
\end{enunciadoLiteral}

\ifgabarito
    % --- CAIXA LARANJA (Resolução e Tutor) ---
    \begin{tcolorbox}[colback=CorTitUm!5, colframe=CorTitUm, boxrule=1pt, arc=4pt, left=6pt, right=6pt, top=6pt, bottom=6pt]
    \small\textbf{\textcolor{CorTitUm}{Roteiro do Professor e Tutor IA:}}\par\vspace{0.1cm}
    \color{CinzaEscuro}
    \textbf{1. Ancoragem (Abordagem Socrática):} Conte os palitos da Figura 1. Deu 4? Ok. Agora conte os da Figura 2. Deu 8? Não! Por que deu 7? O que aconteceu com a "parede" que encostou no outro quadrado?\par\vspace{0.1cm}
    
    \textbf{2. Micro-passos da Resolução:}\\
    \textit{Passo 1:} Contagem inicial $\rightarrow$ Fig.1 = 4 palitos, Fig.2 = 7 palitos, Fig.3 = 10 palitos.\\
    \textit{Passo 2:} Identificar a regularidade (razão de crescimento) $\rightarrow$ Do 4 para o 7, aumentou 3. Do 7 para o 10, aumentou 3. Logo, cada novo quadrado acrescenta apenas \textbf{3 palitos} novos (pois compartilha uma haste central).\\
    \textit{Passo 3:} Projetar os próximos $\rightarrow$ Fig.4 seria $10 + 3 = 13$ palitos.\\
    \textit{Passo 4:} Responder a Fig.5 $\rightarrow$ Fig.5 será $13 + 3 = 16$ palitos.\par\vspace{0.1cm}
    
    \textbf{3. Mapeamento de Armadilhas (Alerta Tutor):} O aluno lógico-dedutivo precipitado vai pensar: "Se a Figura 1 tem 4 palitos, a Figura 5 vai ter $5 \cdot 4 = 20$ palitos". O tutor deve exigir que ele desenhe a figura se errar, para ele perceber que os quadrados do meio não usam 4 palitos cada, pois dividem a parede vizinha.
    \end{tcolorbox}
    \vspace{0.15cm}
    
    % --- CAIXA VERDE (Porto Seguro Didático do Professor) ---
    \begin{boxexplicacao}
    \textbf{Porto Seguro Didático (O "Pulo do Gato"):} Sequências são os alicerces da álgebra avançada (PA e PB no Ensino Médio). Ensine os alunos a criarem uma tabelinha no canto do caderno: [Figura - Palitos]. Olhando para os números enfileirados (4, 7, 10, ...), fica óbvio para qualquer criança que a "tabuada oculta" ali é a de contar de 3 em 3.
    \end{boxexplicacao}
\fi

\par\vspace{\espacoQuestao}
\noindent\textcolor{CorLinha}{\rule{\linewidth}{0.5pt}} 
\par\vspace{\espacoPreQuestao}
\subsubsection*{25.} 
\begin{enunciadoLiteral}
Uma professora de Matemática escreveu a seguinte sequência de números na lousa, mas apagou o quinto termo de propósito, deixando um espaço representado pela letra $y$:

$ 5, \ 11, \ 17, \ 23, \ y, \ 35, \ 41 $

Determine o valor numérico que deve estar no lugar de $y$ para que a sequência mantenha a sua regularidade e explique qual é a regra de formação (o padrão) utilizada para ir de um número para o próximo.
\end{enunciadoLiteral}

\ifgabarito
    % --- CAIXA LARANJA (Resolução e Tutor) ---
    \begin{tcolorbox}[colback=CorTitUm!5, colframe=CorTitUm, boxrule=1pt, arc=4pt, left=6pt, right=6pt, top=6pt, bottom=6pt]
    \small\textbf{\textcolor{CorTitUm}{Roteiro do Professor e Tutor IA:}}\par\vspace{0.1cm}
    \color{CinzaEscuro}
    \textbf{1. Ancoragem (Abordagem Socrática):} Como você faz para descobrir a distância (a diferença) entre o 5 e o 11? Essa mesma distância funciona do 11 para o 17?\par\vspace{0.1cm}
    
    \textbf{2. Micro-passos da Resolução:}\\
    \textit{Passo 1:} Analisar a diferença entre os primeiros termos $\rightarrow$ $11 - 5 = 6$. Aumentou 6.\\
    \textit{Passo 2:} Testar se o padrão se repete $\rightarrow$ $11 + 6 = 17$ (Certo). $17 + 6 = 23$ (Certo).\\
    \textit{Passo 3:} Definir a regra de formação $\rightarrow$ A sequência aumenta de 6 em 6 (adiciona-se 6 ao termo anterior).\\
    \textit{Passo 4:} Calcular o valor de $y$ $\rightarrow$ $ 23 + 6 = 29 $. Logo, $y = 29$.\\
    \textit{Passo 5:} Tirar a prova real com o próximo $\rightarrow$ $29 + 6 = 35$ (Bate com a sequência da lousa).\par\vspace{0.1cm}
    
    \textbf{3. Mapeamento de Armadilhas (Alerta Tutor):} Alunos podem tentar multiplicar ou achar que a regra muda no meio do caminho. Exija sempre a "prova real", testando se o número descoberto ($29$) se conecta corretamente com o próximo número já revelado ($35$).
    \end{tcolorbox}
    \vspace{0.15cm}
    
    % --- CAIXA VERDE (Porto Seguro Didático do Professor) ---
    \begin{boxexplicacao}
    \textbf{Porto Seguro Didático (O "Pulo do Gato"):} Brinque de "Detetive de Padrões". O detetive sempre olha as pegadas (os números). Se o passo for sempre do mesmo tamanho (+6), achamos a regra. Isso é a base absoluta das Progressões Aritméticas que eles verão no Ensino Médio, mas ensinada de forma intuitiva.
    \end{boxexplicacao}
\fi

\par\vspace{\espacoQuestao}

\noindent\textcolor{CorLinha}{\rule{\linewidth}{0.5pt}} 
\par\vspace{\espacoPreQuestao}
\subsubsection*{26.} 
\begin{enunciadoLiteral}
Para uma festa na escola, as mesas do refeitório (representadas por retângulos escuros) estão sendo unidas em fileiras. As cadeiras (representadas por círculos claros) são colocadas ao redor das mesas. 

Observe como o número de cadeiras aumenta conforme juntamos mais mesas:

\begin{center}
% ==========================================
% CONTROLE INDIVIDUAL DESTA IMAGEM:
% ==========================================
\begin{adjustbox}{trim=0cm 0cm 0cm 0cm, clip, max width=0.85\linewidth}
\begin{tikzpicture}
    % --- APLICANDO DROP SHADOW E REGRAS PREMIUM ---
    \definecolor{mesaCor}{HTML}{8E44AD}
    \definecolor{cadeiraCor}{HTML}{F1C40F}
    \tikzset{
        mesa/.style={fill=mesaCor, rounded corners=2pt, drop shadow={opacity=0.3, shadow xshift=2pt, shadow yshift=-2pt}},
        cadeira/.style={circle, fill=cadeiraCor, draw=white, thick, inner sep=4pt, drop shadow={opacity=0.3, shadow xshift=1pt, shadow yshift=-1pt}}
    }
    
    % CONFIG 1: 1 Mesa
    \node[font=\bfseries, text=black] at (1, 2) {1 Mesa};
    \draw[mesa] (0.2, 0.5) rectangle (1.8, 1.3);
    % Cadeiras
    \node[cadeira] at (0, 0.9) {}; % esq
    \node[cadeira] at (2, 0.9) {}; % dir
    \node[cadeira] at (1, 1.5) {}; % cima
    \node[cadeira] at (1, 0.3) {}; % baixo
    
    % CONFIG 2: 2 Mesas
    \node[font=\bfseries, text=black] at (4.5, 2) {2 Mesas juntas};
    \draw[mesa] (3, 0.5) rectangle (6, 1.3);
    % Cadeiras
    \node[cadeira] at (2.8, 0.9) {}; % esq
    \node[cadeira] at (6.2, 0.9) {}; % dir
    \node[cadeira] at (3.8, 1.5) {}; % cima 1
    \node[cadeira] at (5.2, 1.5) {}; % cima 2
    \node[cadeira] at (3.8, 0.3) {}; % baixo 1
    \node[cadeira] at (5.2, 0.3) {}; % baixo 2
    
    % CONFIG 3: 3 Mesas
    \node[font=\bfseries, text=black] at (9, 2) {3 Mesas juntas};
    \draw[mesa] (7, 0.5) rectangle (11, 1.3);
    % Cadeiras
    \node[cadeira] at (6.8, 0.9) {}; % esq
    \node[cadeira] at (11.2, 0.9) {}; % dir
    \node[cadeira] at (7.8, 1.5) {}; % cima 1
    \node[cadeira] at (9.0, 1.5) {}; % cima 2
    \node[cadeira] at (10.2, 1.5) {}; % cima 3
    \node[cadeira] at (7.8, 0.3) {}; % baixo 1
    \node[cadeira] at (9.0, 0.3) {}; % baixo 2
    \node[cadeira] at (10.2, 0.3) {}; % baixo 3

\end{tikzpicture}
\end{adjustbox}
\end{center}

Se juntarmos \textbf{6 mesas} em uma única fileira, quantas cadeiras caberão ao redor dessa fileira? Explique a sua lógica de contagem sem precisar desenhar todas as mesas.
\end{enunciadoLiteral}

\ifgabarito
    % --- CAIXA LARANJA (Resolução e Tutor) ---
    \begin{tcolorbox}[colback=CorTitUm!5, colframe=CorTitUm, boxrule=1pt, arc=4pt, left=6pt, right=6pt, top=6pt, bottom=6pt]
    \small\textbf{\textcolor{CorTitUm}{Roteiro do Professor e Tutor IA:}}\par\vspace{0.1cm}
    \color{CinzaEscuro}
    \textbf{1. Ancoragem (Abordagem Socrática):} Toda vez que a gente gruda uma mesa nova na ponta, a gente perde o espaço da cadeira que ficava no meio. Olhando a imagem, cada mesa nova traz quantas cadeiras a mais para a festa?\par\vspace{0.1cm}
    
    \textbf{2. Micro-passos da Resolução:}\\
    \textit{Passo 1:} Contar o estado inicial $\rightarrow$ 1 mesa = 4 cadeiras.\\
    \textit{Passo 2:} Contar os próximos e achar a razão $\rightarrow$ 2 mesas = 6 cadeiras. 3 mesas = 8 cadeiras. Aumenta de 2 em 2 (uma cadeira em cima, outra embaixo da mesa nova).\\
    \textit{Passo 3:} Projetar para 6 mesas $\rightarrow$ 4 mesas = 10 cadeiras; 5 mesas = 12 cadeiras; 6 mesas = 14 cadeiras.\\
    \textit{Passo 4:} Lógica Algébrica (Alternativa) $\rightarrow$ A regra é: o dobro do número de mesas mais as 2 cadeiras das pontas ($2n + 2$). Para 6 mesas: $2 \cdot 6 + 2 = 12 + 2 = 14$.\par\vspace{0.1cm}
    
    \textbf{3. Mapeamento de Armadilhas (Alerta Tutor):} O erro de proporcionalidade direta é fatal aqui. O aluno pensa: "Se 3 mesas têm 8 cadeiras, 6 mesas terão o dobro (16)". Refute essa lógica desenhando que a junção "rouba" as cadeiras que ficariam entre as mesas.
    \end{tcolorbox}
    \vspace{0.15cm}
    
    % --- CAIXA VERDE (Porto Seguro Didático do Professor) ---
    \begin{boxexplicacao}
    \textbf{Porto Seguro Didático (O "Pulo do Gato"):} Esse é o melhor exemplo visual para mostrar que a álgebra modela a vida real. Mostre que não importa quantas mesas colocarmos, sempre teremos as cadeiras "de cima e de baixo" ($2 \cdot n$) e as famosas cadeiras dos "chefes da mesa" nas duas pontas ($+2$). Regra visual pura!
    \end{boxexplicacao}
\fi

\par\vspace{\espacoQuestao}

\noindent\textcolor{CorLinha}{\rule{\linewidth}{0.5pt}}
\par\vspace{\espacoPreQuestao}
\subsubsection*{27.} 
\begin{enunciadoLiteral}
Um grupo de estudos estava analisando o crescimento das visualizações de um vídeo na internet. No primeiro dia, houve \textbf{2 visualizações}. No segundo dia, foram \textbf{4}. No terceiro dia, \textbf{8}. No quarto dia, \textbf{16 visualizações}.

Um dos alunos anotou no quadro: \textit{"A regra da sequência 2, 4, 8, 16 é que estamos sempre somando 2 ao número anterior."}

\begin{boxresponda}
\textbf{Responda:} A afirmação do aluno está correta? Caso esteja errada, explique o que realmente está acontecendo matematicamente com os números dia após dia e qual será o número de visualizações no quinto dia.
\end{boxresponda}
\end{enunciadoLiteral}

\ifgabarito
    % --- CAIXA LARANJA (Resolução e Tutor) ---
    \begin{tcolorbox}[colback=CorTitUm!5, colframe=CorTitUm, boxrule=1pt, arc=4pt, left=6pt, right=6pt, top=6pt, bottom=6pt]
    \small\textbf{\textcolor{CorTitUm}{Roteiro do Professor e Tutor IA:}}\par\vspace{0.1cm}
    \color{CinzaEscuro}
    \textbf{1. Ancoragem (Abordagem Socrática):} Se fosse verdade que está sempre somando 2, depois do 4 viria qual número? E na nossa sequência, veio o 8. Qual operação transforma o 4 em 8 e o 8 em 16?\par\vspace{0.1cm}
    
    \textbf{2. Micro-passos da Resolução:}\\
    \textit{Passo 1:} Diagnosticar o erro $\rightarrow$ A afirmação está errada. Se estivesse somando 2, a sequência seria 2, 4, 6, 8, 10.\\
    \textit{Passo 2:} Identificar a regra real $\rightarrow$ A sequência está dobrando de valor a cada dia (multiplicando por 2).\\
    \textit{Passo 3:} Calcular o quinto dia $\rightarrow$ Se no quarto dia foram 16, no quinto será o dobro de 16.\\
    \textit{Passo 4:} Finalizar a conta $\rightarrow$ $ 16 \cdot 2 = 32 $ visualizações.\par\vspace{0.1cm}
    
    \textbf{3. Mapeamento de Armadilhas (Alerta Tutor):} Crianças frequentemente testam a regra apenas nos dois primeiros números. De 2 para 4, somar 2 funciona, então o cérebro "aceita" a regra. O tutor deve ensinar o hábito de testar a hipótese no terceiro termo antes de validá-la.
    \end{tcolorbox}
    \vspace{0.15cm}
    
    % --- CAIXA VERDE (Nota Pedagógica) ---
    \begin{boxexplicacao}
    \textbf{Justificativa Curricular (Raio-X da Questão):} Introduz sutilmente o comportamento multiplicativo (exponencial e PG) em oposição ao aditivo (linear e PA), treinando a flexibilidade cognitiva para não achar que toda sequência é baseada em soma.
    \end{boxexplicacao}
\fi

\par\vspace{\espacoQuestao}

\noindent\textcolor{CorLinha}{\rule{\linewidth}{0.5pt}} 
\par\vspace{\espacoPreQuestao}
\subsubsection*{28.} 
\begin{enunciadoLiteral}
Lucas começou a poupar dinheiro em seu cofrinho. O gráfico em barras abaixo mostra a quantia acumulada por ele, em reais, ao longo dos primeiros três meses.

\begin{center}
% ==========================================
% CONTROLE INDIVIDUAL DESTA IMAGEM:
% ==========================================
\begin{adjustbox}{trim=0cm 0cm 0cm 0cm, clip, max width=0.6\linewidth}
\begin{tikzpicture}
    % --- APLICANDO DROP SHADOW E REGRAS PREMIUM ---
    \definecolor{barraCor}{HTML}{2ECC71}
    \definecolor{eixoCor}{HTML}{2C3E50}
    
    \begin{axis}[
        ybar,
        bar width=20pt,
        axis lines=left,
        axis line style={eixoCor, thick},
        enlarge x limits=0.25,
        ymin=0, ymax=60,
        ylabel={\textbf{Valor (R\$)}},
        xlabel={\textbf{Meses}},
        xtick=data,
        xticklabels={Mês 1, Mês 2, Mês 3},
        nodes near coords,
        nodes near coords align={vertical},
        every node near coord/.append style={font=\bfseries},
        tick label style={font=\small\bfseries},
        label style={font=\bfseries},
        hide x axis=false,
        axis x line=bottom,
        axis y line=left,
        ymajorgrids=true,
        grid style={dashed, gray!30}
    ]
    % Adicionando sombra manualmente via estilo não nativo do pgfplots é complexo, 
    % vamos usar filldenso e borda arredondada na definição do plot
    \addplot[fill=barraCor, draw=eixoCor, thick, rounded corners=1pt] coordinates {
        (1, 20)
        (2, 32)
        (3, 44)
    };
    \end{axis}
\end{tikzpicture}
\end{adjustbox}
\end{center}

Se Lucas mantiver exatamente o mesmo padrão de economia todos os meses (poupando o mesmo valor fixo mensalmente), qual será a quantia total, em reais, que ele terá acumulado ao chegar no \textbf{Mês 6}?
\end{enunciadoLiteral}

\ifgabarito
    % --- CAIXA LARANJA (Resolução e Tutor) ---
    \begin{tcolorbox}[colback=CorTitUm!5, colframe=CorTitUm, boxrule=1pt, arc=4pt, left=6pt, right=6pt, top=6pt, bottom=6pt]
    \small\textbf{\textcolor{CorTitUm}{Roteiro do Professor e Tutor IA:}}\par\vspace{0.1cm}
    \color{CinzaEscuro}
    \textbf{1. Ancoragem (Abordagem Socrática):} Se ele tinha 20 e foi para 32, quantos reais ele colocou no cofrinho no segundo mês? Será que no terceiro mês ele colocou esse mesmo valor?\par\vspace{0.1cm}
    
    \textbf{2. Micro-passos da Resolução:}\\
    \textit{Passo 1:} Ler o gráfico e identificar a sequência $\rightarrow$ Mês 1: 20, Mês 2: 32, Mês 3: 44.\\
    \textit{Passo 2:} Achar a razão de crescimento $\rightarrow$ $32 - 20 = 12$. A cada mês, ele deposita 12 reais.\\
    \textit{Passo 3:} Projetar Mês 4 $\rightarrow$ $44 + 12 = 56$.\\
    \textit{Passo 4:} Projetar Mês 5 $\rightarrow$ $56 + 12 = 68$.\\
    \textit{Passo 5:} Projetar Mês 6 $\rightarrow$ $68 + 12 = 80$.\\
    \textit{Passo 6:} Resposta $\rightarrow$ Lucas terá $ \text{R\$} \ 80,00 $.\par\vspace{0.1cm}
    
    \textbf{3. Mapeamento de Armadilhas (Alerta Tutor):} Cuidado com a interpretação de que o valor do mês é $12 \cdot 6 = 72$. A regra não é "ele guarda 12", a regra é "ele JÁ TINHA um começo". Ele começou com 20, então o 20 faz parte da base. Se o aluno errar, faça-o somar passo a passo.
    \end{tcolorbox}
    \vspace{0.15cm}
    
    % --- CAIXA VERDE (Porto Seguro Didático do Professor) ---
    \begin{boxexplicacao}
    \textbf{Porto Seguro Didático (O "Pulo do Gato"):} Gráficos de barras que crescem em "escadinha" igual mostram perfeitamente o que é uma regularidade linear. O tamanho do "degrau" que você sobe de uma barra para a outra é sempre o mesmo. Neste caso, o degrau vale exatos 12 reais.
    \end{boxexplicacao}
\fi

\par\vspace{\espacoQuestao}

\noindent\textcolor{CorLinha}{\rule{\linewidth}{0.5pt}} 
\par\vspace{\espacoPreQuestao}
\subsubsection*{29.} 
\begin{enunciadoLiteral}
Uma sequência numérica é dada pelos múltiplos não nulos do número 5, ou seja: $ 5, 10, 15, 20, 25 \dots $. A expressão algébrica que define qualquer termo dessa sequência é $5n$, onde $n$ é a posição do número (posição 1, posição 2, etc).

Seguindo esse mesmo raciocínio, qual das alternativas abaixo apresenta a expressão algébrica correta para gerar a sequência $ 4, 8, 12, 16, 20 \dots $ ? 

Justifique por que as outras três expressões estão erradas, testando-as para a primeira posição ($n=1$).

\begin{enumerate}[label=\alph*), itemsep=\espacoAlt]
    \item $ n + 4 $
    \item $ 4n $
    \item $ n \cdot n $ (ou $n^2$)
    \item $ 2n + 2 $
\end{enumerate}
\end{enunciadoLiteral}

\ifgabarito
    % --- CAIXA LARANJA (Resolução e Tutor) ---
    \begin{tcolorbox}[colback=CorTitUm!5, colframe=CorTitUm, boxrule=1pt, arc=4pt, left=6pt, right=6pt, top=6pt, bottom=6pt]
    \small\textbf{\textcolor{CorTitUm}{Roteiro do Professor e Tutor IA:}}\par\vspace{0.1cm}
    \color{CinzaEscuro}
    \textbf{1. Ancoragem (Abordagem Socrática):} Se a tabuada do 5 usa a fórmula $5n$, como será a fórmula da tabuada do 4?\par\vspace{0.1cm}
    
    \textbf{2. Micro-passos da Resolução:}\\
    \textit{Passo 1:} Validar a certa (B) $\rightarrow$ A sequência é a tabuada do 4. Logo, é multiplicar a posição por 4 ($4n$). Teste para $n=1$: $4 \cdot 1 = 4$ (Certo).\\
    \textit{Passo 2:} Refutar A ($n+4$) $\rightarrow$ Testando para a posição 1 ($n=1$), teremos $1+4=5$. O primeiro número seria 5 e não 4.\\
    \textit{Passo 3:} Refutar C ($n^2$) $\rightarrow$ Testando para $n=1$, teremos $1 \cdot 1 = 1$. O primeiro número seria 1.\\
    \textit{Passo 4:} Refutar D ($2n+2$) $\rightarrow$ Testando para $n=1$, teremos $2 \cdot 1 + 2 = 4$. Deu certo no primeiro! Mas vamos testar $n=2$: $2 \cdot 2 + 2 = 6$. O segundo número seria 6, e na sequência o segundo é 8. Logo, é falsa.\par\vspace{0.1cm}
    
    \textbf{3. Mapeamento de Armadilhas (Alerta Tutor):} A alternativa D é a clássica "pegadinha" matemática. Ela funciona para o primeiro termo, criando uma falsa sensação de segurança. O aluno OBRIGATORIAMENTE precisa testar o segundo termo para refutá-la.
    \end{tcolorbox}
    \vspace{0.15cm}
    
    % --- CAIXA VERDE (Porto Seguro Didático do Professor) ---
    \begin{boxexplicacao}
    \textbf{Porto Seguro Didático (O "Pulo do Gato"):} Substituir a posição ($n$) na fórmula é o mesmo que perguntar para o oráculo matemático: "Qual é a senha do cofre número 1?". Se o oráculo der a senha 5 (como na alternativa A), você sabe que a fórmula está quebrada, pois o cofre 1 deve ter a senha 4.
    \end{boxexplicacao}
\fi

\par\vspace{\espacoQuestao}

\noindent\textcolor{CorLinha}{\rule{\linewidth}{0.5pt}} 
\par\vspace{\espacoPreQuestao}
\subsubsection*{30.} 
\begin{enunciadoLiteral}
(Desafio Visual) Um arquiteto está projetando a paginação de um piso usando ladrilhos claros e escuros. Ele criou um padrão de "ilhas", onde o ladrilho escuro central (representando $n$) é sempre rodeado por uma moldura de ladrilhos claros, como mostra o projeto:

\begin{center}
% ==========================================
% CONTROLE INDIVIDUAL DESTA IMAGEM:
% ==========================================
\begin{adjustbox}{trim=0cm 0cm 0cm 0cm, clip, max width=0.85\linewidth}
\begin{tikzpicture}
    % --- APLICANDO DROP SHADOW E REGRAS PREMIUM ---
    \definecolor{escuro}{HTML}{2C3E50}
    \definecolor{claro}{HTML}{ECF0F1}
    \definecolor{rejunte}{HTML}{95A5A6}
    
    \tikzset{
        blocoEscuro/.style={fill=escuro, draw=rejunte, thick, drop shadow={opacity=0.3, shadow xshift=1pt, shadow yshift=-1pt}},
        blocoClaro/.style={fill=claro, draw=rejunte, thick, drop shadow={opacity=0.3, shadow xshift=1pt, shadow yshift=-1pt}}
    }
    
    % FIGURA 1 (1 escuro, 8 claros)
    \node[font=\bfseries, text=black] at (1, 3.5) {Projeto 1};
    % claros ao redor
    \draw[blocoClaro] (0, 0) rectangle (1, 1);
    \draw[blocoClaro] (1, 0) rectangle (2, 1);
    \draw[blocoClaro] (2, 0) rectangle (3, 1);
    \draw[blocoClaro] (0, 1) rectangle (1, 2);
    \draw[blocoClaro] (2, 1) rectangle (3, 2);
    \draw[blocoClaro] (0, 2) rectangle (1, 3);
    \draw[blocoClaro] (1, 2) rectangle (2, 3);
    \draw[blocoClaro] (2, 2) rectangle (3, 3);
    % escuro centro
    \draw[blocoEscuro] (1, 1) rectangle (2, 2);
    
    % FIGURA 2 (2 escuros, 10 claros)
    \node[font=\bfseries, text=black] at (6, 3.5) {Projeto 2};
    \draw[blocoClaro] (4, 0) rectangle (5, 1);
    \draw[blocoClaro] (5, 0) rectangle (6, 1);
    \draw[blocoClaro] (6, 0) rectangle (7, 1);
    \draw[blocoClaro] (7, 0) rectangle (8, 1);
    
    \draw[blocoClaro] (4, 1) rectangle (5, 2);
    \draw[blocoClaro] (7, 1) rectangle (8, 2);
    
    \draw[blocoClaro] (4, 2) rectangle (5, 3);
    \draw[blocoClaro] (5, 2) rectangle (6, 3);
    \draw[blocoClaro] (6, 2) rectangle (7, 3);
    \draw[blocoClaro] (7, 2) rectangle (8, 3);
    % escuros
    \draw[blocoEscuro] (5, 1) rectangle (6, 2);
    \draw[blocoEscuro] (6, 1) rectangle (7, 2);

    % FIGURA 3 (3 escuros, 12 claros)
    \node[font=\bfseries, text=black] at (11.5, 3.5) {Projeto 3};
    \foreach \x in {9, 10, 11, 12, 13} {
        \draw[blocoClaro] (\x, 0) rectangle (\x+1, 1);
        \draw[blocoClaro] (\x, 2) rectangle (\x+1, 3);
    }
    \draw[blocoClaro] (9, 1) rectangle (10, 2);
    \draw[blocoClaro] (13, 1) rectangle (14, 2);
    % escuros
    \draw[blocoEscuro] (10, 1) rectangle (11, 2);
    \draw[blocoEscuro] (11, 1) rectangle (12, 2);
    \draw[blocoEscuro] (12, 1) rectangle (13, 2);

\end{tikzpicture}
\end{adjustbox}
\end{center}

Observe bem a quantidade de ladrilhos claros em relação ao número de ladrilhos escuros em cada projeto. Se o arquiteto fizer um projeto gigantesco utilizando \textbf{20 ladrilhos escuros} no centro, quantos ladrilhos claros ele precisará comprar para fazer toda a moldura?
\end{enunciadoLiteral}

\ifgabarito
    % --- CAIXA LARANJA (Resolução e Tutor) ---
    \begin{tcolorbox}[colback=CorTitUm!5, colframe=CorTitUm, boxrule=1pt, arc=4pt, left=6pt, right=6pt, top=6pt, bottom=6pt]
    \small\textbf{\textcolor{CorTitUm}{Roteiro do Professor e Tutor IA:}}\par\vspace{0.1cm}
    \color{CinzaEscuro}
    \textbf{1. Ancoragem (Abordagem Socrática):} Desenhe só a linha de cima e a de baixo de ladrilhos claros. Elas têm a mesma quantidade, certo? E além delas, quantos ladrilhos ficam fixos fechando as duas laterais nas pontas?\par\vspace{0.1cm}
    
    \textbf{2. Micro-passos da Resolução:}\\
    \textit{Passo 1:} Contar os padrões $\rightarrow$ Proj.1: 1 escuro = 8 claros. Proj.2: 2 escuros = 10 claros. Proj.3: 3 escuros = 12 claros.\\
    \textit{Passo 2:} Extrair a lógica visual $\rightarrow$ Para cada ladrilho escuro $n$, existe um claro em cima dele e um embaixo ($2 \cdot n$). Além disso, sempre existem 6 claros fixos nas laterais (3 de cada lado). Ou seja: $ 2n + 6 $.\\
    \textit{Passo 3:} Calcular para $n = 20$ $\rightarrow$ Se são 20 escuros, teremos 20 claros em cima e 20 claros embaixo (total 40). Mais os 6 das pontas laterais.\\
    \textit{Passo 4:} Expressão Final $\rightarrow$ $ 2 \cdot 20 + 6 = 40 + 6 = 46 $ ladrilhos claros.\par\vspace{0.1cm}
    
    \textbf{3. Mapeamento de Armadilhas (Alerta Tutor):} Sem modelagem algébrica, o aluno tentará desenhar 20 ladrilhos escuros no caderno e se perderá na contagem. Aja como mentor: "E se fossem 1000 ladrilhos, você desenharia? Tente achar a conta que funciona para qualquer tamanho."
    \end{tcolorbox}
    \vspace{0.15cm}
    
    % --- CAIXA VERDE (Porto Seguro Didático do Professor) ---
    \begin{boxexplicacao}
    \textbf{Porto Seguro Didático (O "Pulo do Gato"):} Este é o ápice da introdução à álgebra. O aluno percebe que o desenho ("geometria") e os números ("álgebra") são a mesma coisa falada em línguas diferentes. Separar as laterais constantes ($+6$) das partes que crescem ($2n$) é o melhor exercício mental para preparar a mente deles para Funções.
    \end{boxexplicacao}
\fi

\par\vspace{\espacoQuestao}

\end{multicols*}
\end{document}